\chapter{The EARTHFILES}\label{ch:earthfiles}
\markboth{The Earthfiles}{The Earthfiles}

\chapquote{``Hey, what if those crop circles are just ads for Target?''}{Dennis Miller}

The premise of this chapter derives itself from a weekly podcast of the same name. Edited and
produced by documentary filmmaker and investigative journalist Linda Moulton Howe, Howe's weekly
show deals primarily with --- for lack of better terms --- what can only be described as real-life
X-Files cases.

Like many of the invented cases and stories of that 1990s television show, the real-life
counterparts are often dark and mysterious, and otherwise have no scientific, rational, or sound
justification for their existence. For the completeness of information, I have chosen to include
various details and photographs of these cases. Many of which are not pretty to look at, and very
upsetting to read. Due to the nature of this information, I feel the need to warn you, the
student, that some of what you are about to read is extremely disturbing.

I also want to set out how this chapter is going to work, because it is not the same as the
chapters before it. Every case here is associated, in the public mind and usually in the archival
record as well, with a single investigator. That is not an accident of my sourcing. It is a fact
about the field. When a rancher in Colorado finds something he cannot explain, he does not call a
university pathology department. He calls the person who has been on the radio talking about it.
That pattern --- the private archive, the trusted intermediary, the case file that exists in one
place only --- is itself one of the most important things a student of this subject needs to
understand, and Section~\ref{sec:cattlemutilations} is where we look at it directly.

There is a second reason this chapter sits where it does. We have just come out of two sections
--- Section~\ref{sec:philadelphiaexperiment} and Section~\ref{sec:montauk} --- in which a legend
was assembled out of a real classified programme, a decommissioned base, and the regressed
testimony of a handful of interlinked authors, and in which the eventual output was not a finding
but a television format. Every case in this chapter is a variation on that construction,
performed outdoors and without the bunker. A real and unpleasant biological process becomes a
surgical alien harvest. A real and heavily regulated branch of genetics becomes a hybridisation
programme. Two men with a plank become a message from a higher intelligence. A temperature
inversion becomes the trumpets of Revelation. A slab avalanche becomes a compelling natural
force that nobody is permitted to name. And at the end of the chapter, in
Section~\ref{sec:skinwalker}, all of it arrives simultaneously on one property in Utah, with a
camera crew.

\hi{Read this chapter as a single argument rather than a catalogue.} The individual verdicts
matter less than the recurring shape, and by the last page the shape should be doing the work for
you.

% ---------------------------------------------------------------
\section{Cattle Mutilations}\label{sec:cattlemutilations}

For anyone not already aware of what exactly this is, or why it is so significant, the story often
goes something like this. Farmer Joe wakes up bright eyed and early in the morning to tend to his
farm. Upon arriving inside the barn, he is immediately aware that one of his prize steers is
missing. Nowhere to be found. He takes the tractor out to drive around his 300-acre property and
see if he can find this missing cow. Somewhere along the outskirts of his fence, he finds his
missing cow, but is immediately shocked and horrified by the sight in front of him.

Laying before him, his cow is dead. The cavity of the chest is ripped open with clean cuts, as
though surgical tools were used to commit such a crime. The random gnawing and biting of the skin
from predators such as wolves does not account for the perfectly clean cuts on the belly of the
beast. Inside the cow, Farmer Joe notices that all internal organs are missing. The heart, the
liver, the lungs --- there is nothing but empty space inside.

The face of the cow is rather disfigured, as though all bones had been broken right before it met
its untimely demise. In the sockets where eyes should be, there is nothing; only the frightening
stare into oblivion from this now dead animal remains. The cow has had all fluid from inside
removed. No blood, no water, no urine is to be found. And the grass surrounding the beast is not
touched. Not one drop of blood is to be found, suggesting this final resting location of Farmer
Joe's cow is likely not where the actual killing took place.

The problem is that no sign of dragging the animal to this location can be found. The surrounding
grass does not have any traces of blood, nor has it been flattened in the manner one would expect
with the dragging of a several-hundred-pound animal overtop. It is as though this cow was dropped
there from above, or was magically teleported to that very location.

\figwrap[r]{0.40}{A cow dead of ordinary causes on the Snake River Plain, Idaho. This is not a
mutilation photograph, and I have deliberately not used one: every widely circulated example is
either copyrighted or of unverifiable provenance. What the picture does show is the baseline ---
what an unattended carcass on open range simply looks like.}{https://commons.wikimedia.org/wiki/File:Dead_cow_on_Snake_River_Plain,_Idaho.jpg}{fig:deadcow}

\figwrap[l]{0.38}{A turkey vulture feeding on a washed-up carcass. Scavengers take the softest
tissue first --- eyes, tongue, lips, anus, udder --- which is precisely the list of ``surgical
excisions'' the mutilation literature treats as anomalous.}{https://commons.wikimedia.org/wiki/File:Turkey_Vulture_feeding_on_a_washed_up_carcass.jpg}{fig:vulture}
The local authorities are notified but, upon further investigation, are unable to make heads or
tails of what could possibly account for this vicious and brutal murder. So, in the absence of
answers, Farmer Joe then contacted Linda Moulton Howe --- investigative reporter --- to come take a
look.

\subsection{Linda Moulton Howe and the Making of a Mystery}

\hi{That last sentence is the one to slow down on, because it is the hinge of the entire subject.} The
modern cattle mutilation phenomenon is not primarily a veterinary story or a law-enforcement
story. It is a broadcasting story. And the single person most responsible for its shape, its
vocabulary, and its persistence across five decades is Linda Moulton Howe.

Howe was born in Boise, Idaho, in 1942, and received a master's degree in communication from
Stanford University in 1968. Her early career was conventional, successful, and had nothing to do
with any of this: she worked in local television as a documentary producer, winning regional Emmy
awards and a Peabody for environmental and public-health work, and from 1978 to 1983 she was
director of special projects at KMGH-TV, the CBS affiliate in Denver. Her focus in that period was
primarily environmental --- pollution, water, public health in the Mountain West. She was, in
other words, a working investigative journalist with a portfolio of documented, awarded reporting
before she ever looked at a dead cow. Students of this field should hold on to that, because the
lazy version of the sceptical position is to imply that everyone who investigates anomalies is a
crank who could not get hired elsewhere. Howe could get hired anywhere.

In 1979 and 1980, working out of Denver, she produced and directed \emph{A Strange Harvest}, a
one-hour documentary on the livestock deaths then being reported across eastern Colorado and New
Mexico. It aired in 1980, won a regional Emmy, and is the founding text of the subject. Almost
everything the general public believes about cattle mutilations --- the surgical precision, the
bloodlessness, the cored orifices, the absent tracks, the unmarked helicopters, the reluctance of
officials to engage --- reaches them through that film or through material downstream of it. Howe
interviewed ranchers, sheriff's deputies, and veterinarians on camera, and she reached a
conclusion that she has held ever since: that the excisions were performed with a technology not
available to the ranchers or to conventional predators, and that the most probable explanation
involved a non-human intelligence conducting some form of biological harvest.

She did not stop there, and this is what distinguishes her from the many people who produced one
documentary and moved on. She kept the case files. She published \emph{An Alien Harvest} in 1989
and \emph{Glimpses of Other Realities} in two volumes through the 1990s. She became one of the
most frequent guests on \emph{Coast to Coast AM}, appearing regularly from around 1991 under Art
Bell and continuing under George Noory from 2003 until her run as a regular ended in 2019 ---
Section~\ref{sec:artbell} covers what that platform meant for the field. In 1999 she launched Earthfiles as a
website, and from 2007 onward as the podcast and video channel that gives this chapter its name.
Across that span she has produced something genuinely rare: a continuous, single-authored,
five-decade archive of anomaly reports, with names, dates, locations, and photographs, maintained
by one person who answered the phone when ranchers called.

\begin{funfact}[Why This Matters Methodologically]
Howe's archive is simultaneously the best and the worst thing that happened to cattle mutilation
research. Best, because without it we would have almost no documentation of the 1970s wave at all
--- the sheriff's reports are thin, the necropsies were rarely performed, and the carcasses are
long gone. Worst, because a body of evidence curated by a single investigator who holds a specific
hypothesis is a body of evidence with a selection filter attached, and there is no way to unpick
the filter after the fact. If the interesting cases are the ones that got recorded, and the
interesting cases are the ones that fit, the archive will confirm the hypothesis no matter what
was actually happening in the field. This is not a charge of dishonesty. It is a structural
problem, and it applies to every private archive in this book, including the ones assembled by
people I like.
\end{funfact}

The critical case against Howe's interpretation is not that she invented anything. \hi{It is that her
diagnostic standards were those of a documentary producer rather than a pathologist.} Joe Nickell
and other investigators have described her as credulous --- specifically, as accepting witness
characterisation of wounds as ``surgical'' without independent forensic confirmation, and as
treating the absence of an official explanation as evidence for an extraordinary one. In
\emph{A Strange Harvest} the veterinary voices that support anomaly are foregrounded and the ones
that support scavenging are not. Later Earthfiles material has extended the same method to
claims --- underground bases, government insiders, retrieved biological entities --- where the
sourcing is a single anonymous informant, and where a journalist's normal corroboration
requirements appear to have been suspended.

\subsection{What the Investigations Actually Found}

Beginning in earnest in the mid-1970s and centred on Colorado, New Mexico, Kansas, Nebraska and
Montana, ranchers found dead cattle with what appeared to be surgically precise excisions: an
eye, an ear, the tongue removed at the root, the jaw stripped to clean bone, the rectum and
genitals cored out, and the udder taken. The wounds appeared bloodless. Predators reportedly
avoided the carcasses. There were no tracks. Sometimes there were reports of unmarked black
helicopters.

The scale was real enough to matter: Colorado had hundreds of reports by 1975, the state
governor requested federal help, and in 1979 the FBI opened an investigation of cases in New
Mexico under Senator Harrison Schmitt's pressure. Kenneth Rommel, a retired FBI agent, was hired
to run it on a \$44{,}170 grant from the Law Enforcement Assistance Administration. He spent a
year on it, examined cases as they came in rather than after the trail had gone cold, and in June
1980 --- the same year \emph{A Strange Harvest} aired --- published \emph{Operation Animal
Mutilation}, 297 pages, which remains the most thorough field investigation the subject has ever
received.

The conclusion of that investigation, and of every veterinary pathologist who has examined fresh
cases, is normal death followed by scavenging --- and the mechanism explains every one of the
strange features, which is what a good explanation does.

\textbf{The precision.} Blowfly larvae, ants, and small scavengers begin at the moist,
soft-tissue orifices: eyes, mouth, nostrils, anus, vulva, udder. That is not selective interest in
reproductive organs; it is where the openings are. Larval feeding produces a smooth, scalloped
margin with no tearing, which looks under casual inspection like a cut, and looks under a
microscope like exactly what it is --- pathologists distinguish them without difficulty.

\textbf{The bloodlessness.} Blood settles by gravity and clots within hours; a carcass drained
by livor mortis into the down-side tissues has no free blood at an upper-surface wound. Nothing
was exsanguinated. This is basic post-mortem physiology and it is the single detail most
responsible for the mystery, because ranchers know living animals and pathologists know dead ones.

\textbf{The stripped jaw.} Soft tissue is consumed first; hide contracts and splits along the
line of least resistance as it dries, exposing bone cleanly.

\textbf{The absent tracks.} Hard ground, wind, and the interval between death and discovery.
Rommel documented cases where reported ``no tracks'' sites had tracks visible on examination.

\textbf{The predator avoidance.} Not consistently true, and coyotes are selective about
putrefying meat.

Rommel also found that the reported rate did not exceed baseline cattle mortality. Cattle die
constantly --- of blackleg, anthrax, lightning, bloat, plant toxicity, calving complications and
being struck by lightning while standing under the only tree. What changed in 1975 was not the
death rate. It was that ranchers started looking at carcasses with a hypothesis in mind, and a
media narrative told them what to look for.

\begin{funfact}
The black helicopters were probably real, and that is the most interesting part. The mid-1970s
saw genuine unmarked helicopter activity across the American West, some of it associated with
mineral survey, some with law enforcement narcotics work, and some never accounted for. Also real:
a 1975 wave of livestock deaths coinciding with a documented outbreak of clostridial disease.
Also real: at least some cases were deliberate mutilation by humans, including cult activity and
insurance fraud, and several people were prosecuted. Three real phenomena, one imagined
connection.
\end{funfact}

\subsection{Two Documents, One Year}

So we have a clean natural experiment, and it is worth stating plainly because it is the sort of
thing that only becomes visible when you line the dates up. In 1980, two major works on cattle
mutilations were released. One was a broadcast documentary made by an award-winning journalist who
interviewed the ranchers, filmed the carcasses as she found them, took the witnesses at their word
about what a surgical cut looks like, and concluded that something not of this Earth was
harvesting tissue. The other was a 297-page federally funded field report by a retired FBI agent
who attended fresh scenes, commissioned necropsies, checked the reported ``no tracks'' sites
himself, compared the reported death rate against baseline herd mortality, and concluded
predation.

One of those documents won an Emmy, spawned a genre, and is still cited constantly. The other sits
in a New Mexico state archive. That asymmetry has nothing to do with which one was more careful,
and everything to do with which one was more watchable. If there is a single lesson in this
chapter, it is that the survival of a claim in public memory is determined by its narrative fitness
and not by its evidential quality --- and that a field which gets its case files from broadcasters
rather than from laboratories will inherit the broadcaster's incentives whether anyone intends it
to or not.

I want to be fair to Howe in closing, because it would be easy not to be. She asked questions that
the relevant authorities were, in the 1970s, genuinely not asking. She recorded testimony that
would otherwise have vanished. She has been consistent, public, and reachable for fifty years,
which is more than can be said for most of the anonymous sources in this book. Her interpretation
of the evidence is, I think, wrong, and wrong in a way that a working pathologist could have
corrected in an afternoon. Those two judgements are compatible, and holding both at once is
roughly the whole skill this textbook is trying to teach.

\begin{srcnotes}
\item Rommel, K.\ M.\ (1980). \emph{Operation Animal Mutilation: Report of the District
Attorney, First Judicial District, State of New Mexico}. Federally funded study, LEAA grant
\#79-DF-AX-0057.
\item Howe, L.\ M.\ (1980). \emph{A Strange Harvest} [documentary]. Denver: KMGH-TV.
\item Howe, L.\ M.\ (1989). \emph{An Alien Harvest}. Littleton, CO: Linda Moulton Howe
Productions.
\item Earthfiles.com --- Howe's ongoing case archive and podcast, 1999--present.
\item Nickell, J., various, in \emph{Skeptical Inquirer}, on mutilation claims and investigative
standards.
\end{srcnotes}

% ---------------------------------------------------------------
\section{Animal-Human Hybridization}\label{sec:animalhumanhybridization}

The mutilation literature does not stop at the carcass. It proposes a purpose for the carcass ---
harvested tissue, harvested blood, a breeding or hybridisation programme requiring biological
material in quantity --- and that proposal is testable in a way the mutilations themselves are not.
So we test it. Section~\ref{sec:cattlemutilations} asked whether the cuts required a scalpel;
this section asks whether the thing the scalpel is supposed to be \emph{for} is possible at all.

Chimeras and hybrids are a real, active, and heavily regulated area of biology, and knowing what
is actually possible is the only way to evaluate the claims in Section~\ref{sec:geneticmotivation} and Section~\ref{sec:nephilim}.

The distinctions matter. A \textbf{hybrid} results from fertilisation between two species and
has mixed DNA in every cell --- a mule, a liger. A \textbf{chimera} contains cells from two
organisms in separate lineages, mixed but not blended. A \textbf{transgenic} organism carries
specific inserted genes.

What has actually been done: human--animal chimeras exist in research and are strictly limited.
Human cells have been introduced into pig, sheep and monkey embryos, with human contribution
remaining very low --- typically a tiny fraction of cells --- and the embryos destroyed within
days under protocol. Human--mouse neural chimeras have been made to study brain development.
Human genes are routinely expressed in animals for research and pharmaceutical production, and
human insulin has been made in bacteria since 1978. Pig heart valves have been used in humans for
decades, and in 2022 a genetically modified pig heart was transplanted into a human patient who
survived two months.

What has \emph{not} been done, and cannot be: viable hybridisation between species separated by
any substantial evolutionary distance. The barriers are hard. Chromosome number and structure
must be compatible enough for meiosis --- humans have 46, chimpanzees 48, and even that gap makes
viable fertile offspring effectively impossible. Gamete recognition proteins are
species-specific. Developmental gene regulation is timed differently. Mitochondrial and nuclear
genomes must be co-adapted. Hybrids between close species are usually sterile precisely because
of this: mules cannot breed.

And notice what that implies for the alien hybrid claim. Horse and donkey diverged around four
million years ago, share nearly all their biochemistry, and produce a sterile offspring.
An organism from another biosphere would not share our chirality conventions, our genetic code,
our amino acid set, our nucleotides, or our replication machinery --- if it used nucleic acids at
all. It is not a distant relative. It is not a relative. Section~\ref{sec:geneticmotivation} makes the point at length;
here I will simply say that the abduction breeding programme is not improbable in the way a long
shot is improbable. It is a category error, like proposing to cross a fern with a filing cabinet.

\begin{funfact}
Humans are already chimeras, routinely and without technology. Foetal cells cross the placenta and
persist in the mother for decades --- microchimerism --- and have been found in maternal brain
tissue. Cells can pass between twins in utero. And roughly 8\% of the human genome is
endogenous retroviral sequence: ancient viral DNA integrated into our germ line, some of it now
essential, including a gene co-opted for placental function. We are a mosaic of other things
already. Nature did the interesting version first.
\end{funfact}

\Needspace*{0.40\textheight}
\subsubsection*{The chimera we had already built, and the one we were told about}

\figwrap[r]{0.36}{The Chimera of Arezzo, Etruscan bronze, c.\ 400 BC: lion, goat and serpent in one body. The idea that mixing kinds is an abomination is older than biology by two and a half thousand years, and it arrived in the statute book with its moral force intact and its evidence base never examined.}{https://commons.wikimedia.org/wiki/File:Chimera_d\%27arezzo,_fi,_04.JPG}{fig:chimarezzo}
Here is the thing that ought to bother you about the preceding four paragraphs. Everything in them
is boring. Human cells in pig embryos, destroyed on day twenty-eight. A pig heart in a man in
Baltimore. Insulin from a bacterium since 1978. Node-by-node, the actual science of mixing human
material with animal material is a story of small percentages, short timelines, institutional
review boards and disappointing results. It is one of the least sensational research programmes in
modern medicine and one of the most tightly supervised.

And yet the single most repeated hybridisation claim in ufology is not about any of that. It is
about a room. \hi{The claim is that beneath Archuleta Mesa near Dulce, New Mexico, on the sixth
level of a seven-level joint human--alien facility, there is a corridor of cages known as
``Nightmare Hall'', containing multi-species hybrids --- human--animal crosses kept alive in
various states of assembly and distress.} Level 7, in the same telling, holds thousands of humans
in cold storage vats, labelled and catalogued.

I take that claim apart properly in Section~\ref{sec:dulce}, where it belongs, because the reasons
it fails are mostly not biological: the foundational documents are forgeries that a US Air Force
counterintelligence officer fed to a distressed civilian, the sole eyewitness cannot be shown ever
to have existed, and a facility of that scale cannot be built without leaving spoil, power lines
and a payroll. What I want here is narrower and, I think, more interesting. Suppose you had never
heard of Paul Bennewitz. Suppose you knew nothing about the provenance. Ask only the question this
section is equipped to answer: \emph{could the room exist?}

And the answer is that Nightmare Hall is not merely unevidenced. It is a description of an
achievement well beyond the reach of the entire global life-sciences enterprise, described as
routine, in a corridor, decades ago. We cannot make a viable cross between a horse and a donkey
that can reproduce. Nightmare Hall asks us to believe in a shelf of stable, surviving, actively
maintained crosses between organisms from separate biospheres --- performed by people who,
according to the same literature, also could not keep the paperwork in one filing cabinet. The
claim is not too dark to be true. \hi{It is too competent to be true.}

\figwrap[l]{0.34}{A mule: horse mother, donkey father, four million years of divergence, and
---~almost always~--- sterile. This is the outer edge of what mixing two kinds actually achieves.
Nightmare Hall is alleged to have gone very much further, using organisms from another biosphere,
in a corridor, in the 1970s.}{https://commons.wikimedia.org/wiki/File:Equidae_E._africanus_asinus_x_E._ferus_caballus_(Mule).jpg}{fig:mulehybrid}
Which brings me to the question this whole chapter has been circling, and the reason I think 15.2
is the most important section in the book.

\subsubsection*{Why is science illegal?}

I want to put that question in the crudest possible form, because the polite forms of it have been
asked for fifty years and have changed nothing.

The cattle mutilation literature says that someone is harvesting tissue in secret. The abduction
literature says that someone is running a breeding programme in secret. The Dulce literature says
that someone is keeping hybrids in cages in secret. Every one of those claims is, as far as I can
establish, false. But notice the shape they share, because the shape is not stupid. Every one of
them is a claim that \hi{the real biology is happening somewhere you are not allowed to look, and
that the people doing it are not the people who publish.}

And that shape is not a fantasy. That shape is the actual regulatory architecture of human
biology, and we built it ourselves, in public, on purpose.

Consider what is currently forbidden, and by what mechanism. In the United States, federal money
cannot be used for research in which a human embryo is created for research purposes or destroyed
--- not because Congress passed a law debated on its merits, but because of the
\textbf{Dickey--Wicker Amendment}, a rider first attached to an appropriations bill in 1996 and
re-attached, without independent debate, every single year since. It has never had a standalone
vote. It expires annually and is resurrected annually. An entire field's boundary has been drawn,
for three decades, by a sentence bolted to a spending bill.

Heritable human germline modification is, in effect, prohibited in the United States by the same
device: since the Consolidated Appropriations Act of 2016, the FDA has been forbidden even to
\emph{acknowledge receipt} of an application involving a heritable genetic modification to a human
embryo. Not forbidden to approve one. Forbidden to consider one. There is no route to a yes, and
therefore no route to a supervised, published, criticised attempt --- which means the first person
to do it anywhere in the world was guaranteed to be someone operating outside that system.

\figwrap[r]{0.32}{He Jiankui in Hong Kong, November 2018, days after announcing the birth of the
first gene-edited children. He was jailed for three years. Note what the case actually demonstrates:
the work was not stopped by the ban, it was stopped by publicity, and it happened in the dark
because the dark was the only place it was permitted to happen at all.}{https://commons.wikimedia.org/wiki/File:He_Jiankui_at_Second_International_Summit_on_Human_Genome_Editing_(cropped).jpg}{fig:hejiankui}
And then someone was. In November 2018 \textbf{He Jiankui} announced the birth of twin girls whose
\emph{CCR5} genes he had edited as embryos. He had no meaningful oversight, forged ethical approval,
misrepresented the procedure to the parents, and produced mosaic edits of no clear benefit to
anyone. In December 2019 a Shenzhen court sentenced him to three years in prison. Everything about
the episode was a disgrace. But I would ask you to sit with the mechanism rather than the outrage:
\hi{the prohibition did not prevent the experiment. It only guaranteed that when the experiment
happened, it happened without a protocol, without a review board, without a control, and without a
paper.} That is what a ban buys you in biology. Not absence. Illegibility.

\Needspace*{0.40\textheight}
\subsubsection*{S.1435, and the year the United States tried to outlaw a category}

\figwrap[l]{0.30}{Senator Sam Brownback of Kansas, who introduced the Human--Animal Hybrid
Prohibition Act on 9 July 2009 with twenty co-sponsors. The bill carried a ten-year maximum
sentence, died in committee, and was never voted on.}{https://commons.wikimedia.org/wiki/File:Sam_Brownback_official_portrait_3.jpg}{fig:brownback}
The clearest specimen I know of is a bill that never became law, and I think every ufologist who
has ever repeated the Dulce story should read it, because it is the closest the United States has
yet come to legislating Nightmare Hall out of existence --- and it did so without once asking
whether Nightmare Hall was possible.

On 9 July 2009, Senator Sam Brownback of Kansas introduced \textbf{S.1435}, the Human--Animal
Hybrid Prohibition Act, with twenty co-sponsors. It would have inserted a new §1131 into Title 18
of the US Code --- the criminal code, alongside kidnapping and wire fraud --- creating the offence
of creating, attempting to create, transferring or transporting a human--animal hybrid. The penalty
was up to ten years' imprisonment and a fine of not less than one million dollars, or twice the
amount of any gross pecuniary gain, whichever was greater.

Read the prohibited list carefully, because this is where legislation stops describing biology and
starts describing a bestiary. The bill's definition covered, among other things, a human embryo
into which a non-human cell had been introduced; a hybrid human--animal embryo produced by
fertilising a human egg with non-human sperm, or the reverse; a non-human egg into which a human
somatic cell nucleus had been transferred; a non-human life form engineered such that human gametes
develop inside its body; and --- subparagraph (H) --- \hi{a non-human life form engineered such
that it contains a human brain or a brain derived wholly or predominantly from human neural
tissues.}

That last clause is the tell. It is not written in the language of a discipline. It is written in
the language of a warning. And the findings section says so outright: that such experiments are
``grossly unethical'', that they ``blur the line between human and animal, male and female, parent
and child, and one individual and another'', and that they threaten ``human dignity and the
integrity of the human species''. That is a moral claim of considerable weight. It is also,
notably, not a claim about a risk, a harm, an injury, or a patient. Nothing in the findings
identifies anyone who would be hurt.

S.1435 was read twice and referred to the Judiciary Committee, where it died. It has been
reintroduced in various forms since and has never passed. So its legal effect is nil --- and its
cultural effect is not, because a bill like that is a document in which a legislature writes down
what it is afraid of. Twenty-one senators put their names to the proposition that a category of
organism should be a federal felony carrying the same maximum sentence as trafficking in weapons.
Not a procedure. Not a harm. \hi{A kind of thing that should not exist.}

\subsubsection*{What the bans actually protect}

Here is what I think is going on, and it is the argument I would most like to be remembered for if
this book is remembered for anything.

Biology has, twice, drawn a line well and it is worth naming both, because they are the
counter-example to everything above. In February 1975, 140 scientists, lawyers and journalists met
at \textbf{Asilomar} in California to decide what to do about recombinant DNA. They had proposed
their own moratorium the previous year. They argued for three days, they classified experiments by
risk, they published, and in 1976 the NIH turned their conclusions into the NIH Guidelines --- a
supervision regime, not a prohibition. Recombinant DNA went on to give us insulin, hepatitis B
vaccine, and most of modern pharmacology. Nobody went to prison for it. Second: the fourteen-day
rule on human embryo culture, adopted from the Warnock Report in 1984, was for decades the most
respected line in the field --- and in 2021 the International Society for Stem Cell Research
retired it as a bright line, replacing it with case-by-case review, precisely because it had become
a number rather than a reason.

Both of those are lines drawn by people who could state what harm they were preventing, and who
built a mechanism for changing their minds. Dickey--Wicker cannot state a harm and has no mechanism
for changing its mind; it is renewed by inertia. The FDA germline rider cannot state a harm and
forbids the very process --- application, review, argument --- by which a harm might be
established. S.1435 cannot state a harm and proposed a decade in prison anyway.

So, to answer the question in the heading as directly as I can: \hi{science is illegal wherever
someone found the idea of it disgusting before anyone could establish that it was dangerous.}
Disgust is fast, vivid, and requires no evidence --- and it is therefore enormously better adapted
to the legislative timetable than data is. The result is a boundary drawn in the shape of a taboo,
maintained by procedural devices that nobody has to defend on the record, in a field where the
work does not stop when you forbid it. It only moves.

And that is the sentence that connects this section back to the mutilated cow and the cage in the
mesa. \hi{The Dulce hybrid story is false, but the intuition underneath it --- that the most
consequential biology is happening where the public cannot see it and where the practitioner is not
accountable --- is a fair description of what we have built.} We took the most morally loaded
questions in the life sciences and we removed them from the published record. Then we were
surprised when the public filled the resulting vacuum with corridors of cages.

If you want fewer Nightmare Halls, the remedy is not better debunking. It is fewer unexaminable
prohibitions. A field that can be argued with in public does not need to be imagined.

\begin{funfact}
Dolly the sheep was cloned from a mammary cell in 1996, which is why the Roslin team named her
after Dolly Parton. She required 277 attempts. Within five years, human reproductive cloning had
been banned or restricted in dozens of jurisdictions and condemned by the UN --- a global
prohibition erected in record time against a procedure that, thirty years on, still fails in the
overwhelming majority of attempts in every species tried. Meanwhile the genuinely transformative
descendant of the same work, induced pluripotent stem cells, arrived in 2006 by a completely
different route that no legislature had thought to prohibit, and won Yamanaka the Nobel Prize in
2012. The bans landed on the frightening thing. The important thing came in through a door nobody
was watching.
\end{funfact}

\subsubsection*{The reason the line was drawn: you cannot run a dry run on a person}

\figwrap[r]{0.30}{Dolly, 1996--2003, now in the National Museum of Scotland. Twenty-odd countries
banned human reproductive cloning within five years of her birth. Nobody has ever needed the laws,
because the technique remains hopelessly inefficient --- but everyone got to legislate.}{https://commons.wikimedia.org/wiki/File:Dolly_the_Sheep_National_Museum_of_Scotland.jpg}{fig:dollysheep}

I have spent this section arguing that the prohibitions are badly built. I now have to be fair to
the one of them that had a real argument behind it, because otherwise what follows is worthless.

The germline ban was not only disgust. Underneath it sat a genuine and, in 1996, unanswerable
objection. A heritable edit differs from every other medical intervention in three ways at once: the
patient does not exist yet and cannot consent; the change propagates into descendants who cannot
consent either; and there is no reverting it. Every other procedure in medicine can be stopped,
titrated, or withdrawn. This one ships.

And in 1996 --- and in 2015, when the FDA rider was written --- nobody could tell you in advance what
a given edit would do. That is the part that matters, and it is the part the ufology literature never
notices. The objection was not really about dignity. It was about \hi{prediction}. You could cut a
sequence and you could see whether the cut landed, but you could not say what the resulting protein
would do, what else in the genome carried a similar enough sequence to be cut by accident, or what
the edited allele would cost the organism forty years later. The only instrument that could answer
those questions was the organism. The edit \emph{was} the experiment, and the experiment \emph{was}
a person.

I am a programmer by trade, not a biologist, so let me say it in the only dialect I am actually
qualified in. Germline editing in 1996 was a deployment straight to production, on a live system,
with no staging environment, no test suite, no type checker, no linter, no dry run, and no rollback.
The only way to find out whether the change was correct was to ship it to the one user you could
never roll back. Any engineer would refuse that change, and would be right to. A moratorium on
editing under those conditions was not squeamishness. It was competent engineering practice.

\subsubsection*{A test suite for the genome}

What has changed between then and now is not our courage. It is that the dry run started to exist.

Begin with the state of the ground truth. For most of the era in which these laws were written, the
catalogue of human variants whose effect was actually known was minuscule: a few tens of thousands
of clinically annotated changes against a genome in which roughly \hi{71 million} single-amino-acid
substitutions are chemically possible. That ratio is the whole problem in one number. Nearly every
variant a clinician met was a \emph{variant of uncertain significance} --- which is a formal way of
writing \emph{we have no idea}.

In September 2023, DeepMind published \textbf{AlphaMissense}, which took the protein-structure
machinery of AlphaFold and used it to score every one of those 71 million possible substitutions. It
returned a confident classification for \hi{89\% of them} --- 57\% likely benign, 32\% likely
pathogenic --- and it did so for the entire proteome at once, rather than one gene at a time in one
laboratory at a time. In February 2025 the Arc Institute went a layer deeper with \textbf{Evo 2}, a
foundation model trained on 9.3 trillion nucleotides drawn from around 128,000 genomes spanning every
domain of life, which predicts the effect of mutations \emph{zero-shot} --- with no fine-tuning, on
genes it was never shown labels for --- at better than 90\% accuracy, and, critically, does it in
non-coding DNA as well, which is where the older tools were blind. The same period produced deep
learning models for the other half of the problem: guide-RNA design and off-target prediction, which
is the question of where else in three billion bases your edit might land by accident.

And the clinic caught up faster than the statute books did. \hi{Casgevy}, a CRISPR therapy for
sickle cell disease, was approved by the FDA on 8 December 2023 --- gene editing is now a licensed
medicine with a price. Then, in February 2025, a six-month-old boy named \textbf{KJ Muldoon}, born
with a CPS1 urea-cycle deficiency that kills most of those it does not transplant, received a base
editor \emph{designed for his particular mutation} --- a therapy that existed nowhere on earth
before his diagnosis and was designed, built, tested and dosed in roughly six months. He received
escalating doses between six and eight months of age, tolerated them, and went home. The case was
published in the \emph{New England Journal of Medicine} in June 2025.

So take the two halves together, because separately they are just news. We can now edit a specific
human being's specific mutation on a six-month schedule, and we can now compute, in advance and at
proteome scale, what a given change is likely to do. In the language I know: we finally have a
staging environment and a static analyser. The edit no longer has to be the experiment.

\begin{funfact}[Fun Fact: A Predictor Can Be Wrong. A Rider Cannot.]
\funfactlead{This is the part I would ask a legislator to sit with.}
A variant-effect model publishes a number for every possible mutation, which means it can be
scored, and it \emph{has} been --- against clinical databases, against laboratory mutational scans,
against patient outcomes. When it is wrong you can measure by how much, retrain it, version it, and
compare the new version against the old. It is falsifiable and it improves. The FDA germline rider
forbids the agency even to acknowledge an application, which means it generates no applications, no
reviews, no protocols, no failures and no data. It cannot be scored. It cannot improve. It has been
renewed annually for over a decade without once producing a fact. One of these two instruments is
science and the other is a habit, and the ban is not the one that is science.
\end{funfact}

\subsubsection*{Why now, and what I am not asking for}

Here is the argument, and I will state it as narrowly as I can so that it cannot be mistaken for
something louder.

The germline moratorium was justified by an inability to predict consequences. That inability was
real and it is now partially, measurably, and rapidly less real. When the stated reason for a
prohibition weakens, the prohibition is supposed to be re-examined --- that is the whole difference
between Asilomar and Dickey--Wicker, between a line with a mechanism for changing its mind and a line
without one. Asilomar reconvened, classified experiments by risk, and produced a supervision regime.
The 14-day rule was retired in 2021 when it became a number rather than a reason. The germline ban has
never been re-examined, because the device chosen to implement it --- refusal to acknowledge receipt
--- makes re-examination procedurally impossible. There is no docket. There is nothing to review.

I am not arguing that anyone should edit an embryo tomorrow. I want to be exact about the limits,
because overselling this hands the argument straight back to the people I am arguing with.
Ninety percent accuracy is a research result, not a clinical standard; the eleven percent of variants
AlphaMissense could not call are precisely the interesting ones. These models predict the effect of
variants that \emph{already exist in nature} and have therefore been filtered by selection --- which
is not the same problem as predicting what a novel, deliberate, never-before-observed edit does in a
developing organism. Pleiotropy is unsolved: a change that protects you from one disease may be paying
for it somewhere you are not looking, which is exactly the objection to He Jiankui's \emph{CCR5}
edits and remains valid. Anything polygenic --- height, temperament, intelligence, the entire fantasy
market --- is out of reach and will stay out of reach. And a model trained on the whole tree of life
can tell you what is \emph{unusual}. It cannot tell you what is \emph{good}. That last question was
never going to be answered by a computer.

What I am arguing is smaller and, I think, unanswerable. \hi{The safe way to test this now exists in
silico, and the law forbids the paperwork rather than the experiment.} A regime in which a proposal
must be accompanied by exhaustive computational prediction, adversarial review of that prediction,
demonstration in organoids and animal models, and publication of the failures, is a regime that could
be argued with in public. It is Asilomar with better instruments. The alternative --- the one we
actually have --- is the regime that produced He Jiankui: no route to a yes, therefore no protocol,
no review board, no control, and no paper. We did not prevent the experiment. We exiled it, and then
read about it in the newspaper.

And note, finally, what this does to the story that opened the section. The Dulce hybrid myth imagines
a secret laboratory doing the biology the public is not allowed to see. The genuine version is duller
and worse: the most consequential biology in the species is the one category of work we have made
unexaminable by statute, at exactly the moment we acquired the tools to examine it safely. Which is
the same structural failure this entire book is about --- and the same instrument that fixes it. The
prediction-first, publish-the-failure method I argued for in a lecture hall in 2017 and described in
Section~\ref{sec:ai} is not a ufology technique. It is just how you make a claim that reality is
allowed to refuse. A radar array can do it to a light in the sky. A variant-effect model can do it to
an edit in a genome. In both fields the obstacle was never the absence of a method. It was a
prohibition on being told the answer.

\begin{srcnotes}
\item Cheng, J.\ et al.\ (2023). Accurate proteome-wide missense variant effect prediction with
AlphaMissense. \emph{Science}, 381(6664), eadg7492 --- all 71 million possible missense variants
scored, 89\% classified. \srclink{https://www.science.org/doi/10.1126/science.adg7492}.
\item Brixi, G.\ et al.\ (2025). Genome modelling and design across all domains of life with Evo 2.
Arc Institute / bioRxiv 2025.02.18.638918; 9.3 trillion nucleotides, {\textgreater}90\% zero-shot
variant-effect accuracy. \srclink{https://arcinstitute.org/news/evo2}.
\item US Food and Drug Administration approval of Casgevy (exagamglogene autotemcel) and Lyfgenia for
sickle cell disease, 8 December 2023.
\item Musunuru, K.\ et al.\ (2025). Patient-specific in vivo gene editing to treat a rare genetic
disease. \emph{New England Journal of Medicine}, 392(22), 2235--2243 --- the KJ Muldoon case; first
infusion late February 2025, publicly announced 15 May 2025.
\srclink{https://pubmed.ncbi.nlm.nih.gov/40373211/}.
\item Sherkatghanad, Z.\ et al.\ (2023). Machine learning and deep learning approaches for
CRISPR/Cas9 on-target and off-target prediction --- review. \emph{Briefings in Bioinformatics}.
\item S.1435 --- Human--Animal Hybrid Prohibition Act of 2009, 111th Congress, introduced
9 July 2009 by Sen.\ S.\ Brownback with 20 co-sponsors; read twice and referred to the Committee
on the Judiciary. Full text at \srclink{https://www.congress.gov/bill/111th-congress/senate-bill/1435}.
\item Dickey--Wicker Amendment, first enacted as §128 of P.L.\ 104--99 (1996) and re-enacted as an
annual appropriations rider thereafter.
\item Consolidated Appropriations Act, 2016, P.L.\ 114--113, §749 --- prohibition on FDA
acknowledgement of applications involving heritable modification of the human germline.
\item Berg, P., Baltimore, D., Brenner, S., Roblin, R.\ O., \& Singer, M.\ F.\ (1975). Summary
statement of the Asilomar Conference on Recombinant DNA Molecules. \emph{PNAS}, 72(6), 1981--1984.
\item NIH Guidelines for Research Involving Recombinant DNA Molecules, first issued 1976.
\item International Society for Stem Cell Research (2021). \emph{Guidelines for Stem Cell Research
and Clinical Translation}, revision retiring the 14-day rule as a categorical limit.
\item Wilmut, I.\ et al.\ (1997). Viable offspring derived from fetal and adult mammalian cells.
\emph{Nature}, 385, 810--813.
\item Takahashi, K., \& Yamanaka, S.\ (2006). Induction of pluripotent stem cells from mouse
embryonic and adult fibroblast cultures. \emph{Cell}, 126(4), 663--676.
\item Cyranoski, D., \& Ledford, H.\ (2018--2019). Reporting on the He Jiankui case, \emph{Nature}
news; Shenzhen Nanshan District People's Court judgement, December 2019.
\item Griffin, D.\ E.\ et al.; and Bailey, K.\ V., on xenotransplantation and the 2022 University of
Maryland pig-heart transplant.
\end{srcnotes}

% ---------------------------------------------------------------
\section{Crop Circles}\label{sec:cropcircles}

So far the thread has run on material that resists inspection: a base that was closed, a carcass
that was already decomposing, a breeding programme that exists only as an assertion. This section
is the control experiment. Here the evidence is abundant, the site is accessible, the phenomenon
recurs on a known schedule in a known county --- and, uniquely in this book, the people responsible
put their hands up and demonstrated the method on camera.

This is the cleanest case in the entire book, because we have confessions, we have a
reconstruction, and we can watch the phenomenon respond to human incentives in real time.

Crop formations --- flattened geometric patterns in standing cereal crops --- appeared in numbers
in southern England from the late 1970s, concentrated in Wiltshire and Hampshire around
Avebury, Silbury Hill and the Pewsey Vale. They began as simple circles and became, over two
decades, extraordinarily elaborate: fractals, three-dimensional illusions, and in 1996 the
Julia Set beside Stonehenge, and in 2001 a 400-metre pattern at Milk Hill with 409 circles.

The claimed anomalies were: bent-not-broken stem nodes, altered germination rates, magnetic
particles in soil, electronic equipment failure, dowsing responses, and the sheer complexity.

\figwrap[l]{0.42}{Bower and Chorley's plank-and-rope method, and a formation of the kind declared unfakeable.}{https://commons.wikimedia.org/wiki/File:Crop_circle_Switzerland.jpg}{fig:cropcircle}
In September 1991, Doug Bower and Dave Chorley --- two men in their sixties from Southampton ---
announced that they had made the original formations starting in 1978, and demonstrated the method for the press (see Figure~\ref{fig:cropcircle}): a length of rope, a plank with a loop for the foot, and a baseball cap with
a wire sight for keeping lines straight. They then made a formation and invited a prominent
cereologist to authenticate it, which he did, on camera, declaring that no human could have made
it. That footage exists. Bower and Chorley said they had started in a pub after discussing
Tully, and that they kept going because the reaction was funny.

Since then: circle-making has become an art form with named practitioners, competitions,
commissions from advertising agencies --- Nike, Mitsubishi, Weetabix and others have paid for
them --- and documented overnight construction of complex formations under observation. John
Lundberg's Circlemakers group has built formations for television with cameras running. Physicist
Richard Taylor published in \emph{Nature} in 2011 on the increasing technological sophistication of
circle-making, including the use of GPS and laser levels.

The alleged physical anomalies did not survive testing. Node bending occurs naturally through
phototropic and gravitropic response in flattened stems --- lay a plant down and it bends at the
nodes as it tries to grow upright, within days. Germination differences track mechanical damage.
The magnetic particles are meteoritic dust, which falls everywhere. And the diagnostic test is
decisive: formations appear overnight, in crops, in the specific regions where circle-makers
live, in the months when crops are standing, and they became more complex in step with the
interest and the competition. Crop circles have never appeared in a field under continuous
surveillance in a way that excluded human access, and they follow the aesthetics of human
mathematical fashion --- fractals arrived after fractals were in popular culture.

\begin{funfact}
My favourite detail is the economics. Crop formations cluster in Wiltshire not because of ley
lines but because Wiltshire has tourists, farmers who charge admission, and a concentration of
circle enthusiasts. Some farmers made more from gate money than from the crop. Where there is a
market, there is supply. And the phenomenon has declined markedly since the 2000s --- not because
the visitors left, but because the makers got older and the novelty passed. Anomalies do not
usually retire.
\end{funfact}

\subsubsection*{A portfolio of formations, and what a photograph can and cannot settle}

The four plates that follow are ordinary photographs of ordinary formations, taken by ordinary
photographers who put their names to them. They are reproduced at full page for a reason that cuts
against the way such pictures are normally used. In the cereology literature the aerial photograph
is the argument: the geometry is displayed, the reader is invited to feel that no person could
have walked that out in the dark, and the feeling is treated as evidence. Look at these four at
size and the feeling arrives on schedule. It is worth sitting with it, because the feeling is real
and the inference from it is worthless.

What a photograph of a finished formation can establish is narrow: that the crop was laid flat in
that pattern, at that scale, in that field, before the shutter opened. What it cannot establish is
anything at all about the agency that did the laying. There is no feature of a completed pattern
that distinguishes a plank and a rope from a plasma vortex or a craft, because the only record the
photograph holds is the end state, and the end state is the one thing every competing explanation
agrees on. The causal question is settled by watching the field, not by admiring the result --- and
when fields have been watched, the answer has been men with planks, every time.

So read the plates as the control experiment this section promised. Each one is a formation that
would have been presented in 1990 as beyond human capability, and each one is of a class that
named circle-makers have since built to order, on camera, for money.

\cropplate{pl_crop1}{The formation from the air}{Switzerland, 2009}{%
The standard cereological presentation: altitude, low sun, and a pattern that reads as engineered.
The aerial view is what makes formations persuasive, and it is also what strips out every detail
that would let a viewer judge how the crop was laid. From here the flattening looks seamless. It is
not; see the next plate.}{%
Photograph by Kecko, ``Swiss Crop Circle 2009 Aerial''. Licensed CC BY 2.0 via Wikimedia Commons.}{%
https://commons.wikimedia.org/wiki/File:Swiss_Crop_Circle_2009_Aerial_(3766687033).jpg}

\cropplate{pl_crop2}{The same formation from inside it}{Switzerland, 2009 \mdash{} ground level}{%
Ground level, and the mechanics are visible: stems laid in sweeping bands, a swirl about a centre,
stalks bent at the node rather than snapped. Every one of those was once offered as an anomaly.
Every one of them is what happens when standing cereal is pushed over with a board and then spends
a few days trying to grow upright again. The bands also record the order of work, which is why
makers photograph their own formations from the ground and enthusiasts photograph them from the
air.}{%
Photograph by Kecko, ``Swiss crop circle detail''. Licensed CC BY 2.0 via Wikimedia Commons.}{%
https://commons.wikimedia.org/wiki/File:Swiss_crop_circle_detail.jpg}

\cropplate{pl_crop3}{Geometry at scale}{Diessenhofen, Switzerland, 15 July 2008}{%
A large multi-element design of the kind that arrived after the 1991 confessions, not before. The
developmental sequence is the tell: formations became more complex as the audience, the
competitions and the commissions grew, and the mathematics they quote --- fractals, then
three-dimensional illusions --- followed human intellectual fashion by a few years each time. A
phenomenon tracking a culture's reading habits is a phenomenon with people in it.}{%
Photograph by Hansueli Krapf, ``Aerial View of the Crop Circle in Diessenhofen''. Licensed
CC BY-SA 3.0 via Wikimedia Commons.}{%
https://commons.wikimedia.org/wiki/File:Aerial_View_of_the_Crop_Circle_in_Diessenhofen_15.07.2008_16-44-08.JPG}

\cropplate{pl_crop4}{A field, a footpath, and a queue}{Vacaville, California}{%
The part of the picture the poster prints crop out: the tramlines the tractor left, the
access road along the field edge, and the houses whose windows look straight down onto the
crop. Formations sit beside roads, within walking distance of a gate, in counties with
tourists and farmers willing to charge for parking. That is not a sceptical flourish, it is
a spatial statistic, and it is the same statistic in Wiltshire, in Hampshire and in the
Swiss cantons above.}{%
Photograph by Steve M.~Windham, ``Vacaville CA Crop Circle''. Licensed CC~BY~2.5 via
Wikimedia Commons.}{%
https://commons.wikimedia.org/wiki/File:Vacaville_CA_Crop_Circle.jpg}

\subsubsection*{The Nazca objection, and why old ground markings do not rescue the new ones}

The reply that survives longest, once the plank and the rope are conceded, is a change of subject to
antiquity. If crop formations are made by men in wellingtons, what about the Nazca Lines --- the
hundreds of straight lines, geometric fields and biomorphic figures scored into the desert pavement
of the Peruvian coast between roughly 500~BCE and 500~CE, some of them kilometres long, most of
them legible only from altitude? Erich von D\"aniken made them landing strips. The general form of
the argument is that the ancients could not have built for a viewpoint they could not occupy, so
someone in the air must have commissioned the work; and if that holds for Nazca, ground markings in
general become evidence of visitors, crop formations included.

The archaeology dismantles this at the level of method, and it is worth being specific because the
dismantling is unusually complete. The technique is known: remove the dark, oxidised, iron-stained
surface cobbles and expose the pale gypsum-rich ground a few centimetres beneath. That is trenching
by hand with a basket, and the discarded cobbles are still heaped along the edges of the lines where
the workers left them. Wooden stakes have been recovered at line termini and radiocarbon dated to
the period; Maria Reiche mapped the geometry and showed it is generated by stake-and-cord
construction, the same operation as a builder's string line, extended. Joe Nickell reproduced a
full-size figure in 1982 in a Kentucky field using stakes, cord and a crosspiece, in a matter of
days, with a small crew and no surveying instruments, and the result matched the originals in
precision. The unfinished and mistaken lines are present on the pampa as well --- a corrected
spiral, an abandoned run --- which is exactly what a human workshop leaves and exactly what a
commissioned aerial blueprint does not.

The viewpoint premise is also simply false. The figures are readable from the low hills and the
valley scarps flanking the pampa, from which several of the largest were first recorded in the
modern era, and the makers' interest was in all likelihood not the view at all. The lines converge
on centres and radiate toward water sources and mountain sight lines; many were walked. Anthony
Aveni and Johan Reinhard read them as ritual pathways --- processional routes bound up with water,
fertility and the mountain deities that supply it, in one of the driest inhabited landscapes on
Earth. A line meant to be walked in procession does not need a spectator overhead any more than a
labyrinth on a cathedral floor does.

And the surface argument fails on its own terms. A desert pavement that has held a scratch for two
thousand years, in a place with almost no rain and a stable wind regime, is a preservation
environment, not a runway. Landing anything heavy on it would blow the loose cobbles aside and
destroy the very contrast the figures depend on --- and the pampa shows no compaction, no scorching
and no repeated-use wear where the ``strips'' are. Von D\"aniken's landing field is a surface that
could not survive being landed on once.

So Nazca does not rescue the crop formations; it convicts them by comparison. Both are large ground
markings whose scale outruns a casual observer's intuition about what people will do with rope and
patience. In one case the makers are two thousand years dead and left their tools, their spoil
heaps and their mistakes behind. In the other they are alive, named, filmed, and available for
commission. The pattern is identical in both: the mystery lives entirely in the gap between how
impressive the result looks and how cheap the method turns out to be.

\actualq[assumption]{the geometry is too good for people to have made it}{people are extremely
good at geometry when given rope, stakes, a straight edge and a reason --- and the reason is
usually other people looking}

% ---------------------------------------------------------------
\section{The Trumpets of Jericho}\label{sec:trumpetsjericho}

Section~\ref{sec:cropcircles} showed a phenomenon that scaled with an audience. Keep that in view,
because the case now in front of us is the same economics with the manufacturing cost removed:
instead of a plank, a rope and a wet night in a Wiltshire field, all a contributor needs is an
audio file. And where the crop circles left physical evidence that could be examined, this one
leaves a waveform --- which turns out to be far more incriminating.

Since roughly 2008, and heavily from 2011, people worldwide have recorded low, groaning,
trumpet-like or metallic-grinding sounds apparently coming from the sky, lasting minutes,
sometimes with no visible source. Recordings came from Kiev, Florida, British Columbia, Germany,
Costa Rica, and dozens of other places. They were widely framed as apocalyptic, the trumpets of
Revelation, or the sound of a cloaked craft.

Several distinct real phenomena are in the mix, and the strongest single point about this case is
that the recordings are not all recordings of the same thing.

\textbf{Industrial noise carried by inversion.} A temperature inversion refracts sound downward
and can carry industrial noise --- pile drivers, mining, rail yards, gas flares, ventilation
fans, ice-road construction --- many kilometres, arriving with no apparent direction and
apparently from above. This accounts for a large share, and in several cases the specific source
was identified: a nearby plant, a rail marshalling yard, a snow grooming operation.

\textbf{Ice quakes and glacial calving} produce extraordinary groaning and booming. So does lake
ice --- and the sounds are eerie enough that recordings of them are frequently mislabelled as
sky trumpets.

\textbf{Aurora-related and geomagnetic acoustic effects} are documented but poorly understood.

\textbf{Skyquakes} are a genuine long-standing unexplained category: the Barisal Guns in Bengal,
the Seneca Guns on Seneca Lake and the Carolina coast, \emph{mistpouffers} in Belgium, and
\emph{brontidi} in Italy. These have been reported for centuries and are variously attributed to
supersonic aircraft, distant thunder from over-the-horizon storms, submarine slope failure,
methane release, and minor seismicity. Some remain unexplained and the literature is honest about
it.

\textbf{And an enormous fraction are hoaxes.} This is not speculation. Multiple viral ``sky
trumpet'' videos were traced to the same sound library, and specifically to sound design assets
--- several were identified as an audio track from a film or from the popular ``Red Weather''
stock sample, layered over unrelated footage. Video forensics on the most-shared examples found
identical waveforms across supposedly independent recordings on different continents. \hi{Identical
waveforms are not a coincidence; they are a file.}

\begin{funfact}
The honest position on skyquakes is that a residue of them is unexplained, and I want to put that
plainly because this chapter has been unsympathetic. The Seneca Guns were recorded by James
Fenimore Cooper. Nobody has produced a complete account. But note the shape of the residue:
it is a \emph{sound with no identified source}, not a sound with an identified extraterrestrial
source. Unexplained means unexplained. It does not mean explained by the most exciting hypothesis
available, and if you take one habit from this book, take that one.
\end{funfact}

% ---------------------------------------------------------------
\section{Paranormal Skies}\label{sec:paranormalskies}

That closing point --- that a residue of genuinely unexplained cases exists, and that its existence
licenses nothing in particular --- is the whole subject of this section. The sky trumpets
(Section~\ref{sec:trumpetsjericho}) were a mixed bag of inversions, ice, skyquakes and outright
files. What follows is a longer mixed bag, and I have arranged it deliberately so that you can
watch the sorting happen: most of these have clean mechanisms, a few do not, and the difference is
visible without any special expertise.

\subsubsection*{Light spirals, and why a spiral looks like a wormhole}

I am going to start with the one shape that has done more work in modern ufology than any disc,
because it is the only thing the sky has ever produced that genuinely resembles the textbook
picture of a wormhole: an enormous, silent, geometrically perfect spiral of light, hanging in a
dark sky with a hole in the middle.

\figwrap[r]{0.40}{Twilight phenomenon in the wake of a Minotaur~I launch from Vandenberg, September 2005: unburned propellant freezing in near-vacuum and catching sunlight while the ground below is dark. Give the venting stage a spin and this plume becomes a spiral.}{https://commons.wikimedia.org/wiki/File:Twilight_phenomenon_in_the_wake_of_a_Minotaur_I_launch_from_Vandenberg_AFB,_CaIif.jpg}{fig:spiralplume}
The case that made the shape famous is Norway, 9 December 2009. Shortly before eight in the
morning, in the polar dark, observers across Troms and Finnmark watched a blue-white beam climb
from behind a mountain ridge, stop, and unwind into a vast white spiral of perhaps ten or twelve
turns with a black core, which then faded over several minutes while the beam drooped into a
greenish plume. It was photographed by hundreds of people, including press photographers with
proper equipment, from many separate locations. There is no question about what was seen. It is
probably the best-documented anomalous-looking aerial event in history, and it happened two days
before Barack Obama arrived in Oslo to collect a Nobel Peace Prize, which the internet found
irresistible.

Within a day the Russian Ministry of Defence confirmed a failed test launch of a Bulava
submarine-launched ballistic missile from the White Sea. The third stage malfunctioned. Jonathan
McDowell, the Harvard astrophysicist who tracks launches for a living, laid out the mechanism
immediately: a damaged or asymmetric nozzle imparts a spin to the stage, the stage keeps venting
propellant while it tumbles, and the propellant sprays outward from a rotating source. That is
an Archimedean spiral, and it is what you get every time, because it is geometry rather than
physics. Material leaves the source radially at roughly constant speed while the source turns at
roughly constant rate, so the arms are evenly spaced. In near-vacuum there is no air to smear the
pattern, and the droplets freeze into a fine ice-and-particulate mist that scatters sunlight
brilliantly while the observer on the ground is still in night --- the same twilight-illumination
trick discussed later in this section. The dark hole in the middle is where the venting had
already stopped.

That explanation has a property that almost no ufological explanation has: it makes a measurable
prediction. Count the arms, measure their angular spacing, time the expansion, and you recover the
rotation rate of the object that sprayed them. People did. It matches a tumbling upper stage.
The Norwegian event is not an unexplained case with a proposed solution; it is a solved case,
solved to the level of arithmetic.

\subsubsection*{The spiral is now a monthly event}

The reason I can be confident about Norway is that the shape has since become routine, and it
became routine because of one company's flight profile. When a Falcon~9 second stage performs a
fuel dump or a de-orbit burn at high altitude, and the stage is rolling, it produces the same
spiral. Hawaii in January 2023 gave astronomers at the Subaru telescope a textbook example on
time-lapse; Alaska in April 2023 produced a spiral so clean that the Fairbanks aurora forecasters
had to issue a public explanation; Britain, the Netherlands and northern Europe got one in March
2025 that filled the news cycle for a day. Each time, the launch was on a public schedule, the
timing matched to the minute, and the geometry matched the known roll rate.

This is worth dwelling on because it is a rare gift to this subject. We have a phenomenon that
looks impossible, that people photograph in the tens of thousands, and whose cause is announced
in advance on a website. It is the control group. Anything the sky does that matches this
signature --- silent, expanding, minutes-long, evenly spaced arms, occurring within an hour of
sunrise or sunset, upper atmosphere --- is a venting rocket stage, and you can check the launch
log before you file a report.

Older cases yield to the same key once you know to look. Petrozavodsk, 20 September 1977, is the
Soviet Union's most famous UFO event: a huge jellyfish-shaped luminous body over Karelia, showering
the city with beams of light, reported in \emph{Izvestia} and taken seriously enough to trigger a
formal state investigation that ran for over a decade. It was the launch of Kosmos~955 from
Plesetsk, seen at a distance of a few hundred kilometres, its plume expanding into vacuum and lit
by sunlight the observers on the ground could not see. Chinese launches from Jiuquan and Xichang
have produced regular spiral and jellyfish reports across Asia for the same reason. The category
of ``enormous silent luminous thing in the twilight sky'' is very nearly a synonym for
``somebody's upper stage''.

\subsubsection*{Why a spiral reads as a wormhole}

Now the interesting part, which is psychological rather than physical. Of all the things a spiral
could be mistaken for, why is the immediate popular reading a wormhole or portal? The Norwegian
press said it outright in 2009, and the word has attached itself to every spiral since.

The answer is that we have been trained by pictures. The standard illustration of a wormhole or a
warped spacetime region --- the funnel, the embedding diagram, the whirlpool with a throat at the
centre --- is a spiral seen in perspective, with a dark aperture in the middle. When the sky
produces a luminous spiral with a black core, it matches that trained image far better than it
matches anything in ordinary experience, so that is the label the mind reaches for. It is
pattern-matching to a diagram, not to a thing anybody has ever seen. Nobody has photographed a
wormhole. What people recognise is the textbook figure.

That matters for the case I will come to in Section~\ref{sec:skinwalker}, because a wormhole or
portal is the working hypothesis the Skinwalker Ranch investigators themselves have converged on
for whatever happens over the mesa there --- objects appearing and disappearing, radiation spikes,
and the persistent notion of an opening in the sky that things come through. I want to be scrupulous
about what that hypothesis is worth. It is a guess. It is a guess made by people with instruments
in the field, which is more than most of this subject manages, and it is offered by them as a guess
rather than a finding. But a traversable wormhole requires exotic matter with negative energy
density, of which we have no example, and it has never been observed anywhere by anybody. So the
chain runs: an image from a physics textbook, a shape the atmosphere can make for entirely
pedestrian reasons, and a hypothesis with no confirmed instance in nature. When you hear ``portal''
applied to a light in the sky, that is the chain you are being asked to accept, and the honest
position is to say so and keep watching.

\subsubsection*{Weather, and the working list}

A large fraction of what people see is weather, and weather is stranger than most people's mental
catalogue allows for. Here is the working list I would want anybody to have memorised before
reporting anything.

\textbf{Lenticular clouds} form in standing waves downwind of mountains: smooth, layered,
lens-shaped, stationary while the wind passes through them, often with a hard-edged
disc appearance and stacked ``pile of plates'' structure. They are the single best natural
explanation for a hovering metallic disc, and they hover for hours.

\textbf{Noctilucent clouds}, at around 80 km in the mesosphere, glow electric blue after sunset
long after the sky is dark. \textbf{Nacreous or polar stratospheric clouds} produce iridescent
sheets at 15--25 km.

\textbf{Sun dogs, sun pillars, 22$^\circ$ haloes, circumzenithal arcs and parhelic circles} are
all ice-crystal optics, and they produce geometrically precise bright spots and lines in the sky
that look manufactured because they are governed by crystal geometry.

\textbf{Temperature inversions} bend light, producing superior mirages, looming, and the Fata
Morgana --- which can lift a distant ship, light, or coastline into the sky, stretch it, invert
it, and make it shimmer and appear to move. This is also the principal cause of \emph{radar
anomalies}: ducting in an inversion layer returns ground clutter at apparent altitude and
apparent speed. Section~\ref{sec:uap}'s radar-visual cases need to be read against it.

\textbf{Sprites, blue jets and ELVES} are genuine upper-atmosphere electrical discharges above
thunderstorms, confirmed only in 1989 despite decades of pilot reports that were dismissed ---
which is a case where the aircrew were right and the scientists were wrong, and I include it
because intellectual honesty requires it.

\textbf{Green flashes, mirages of Venus, and atmospheric scintillation} make stars and planets
change colour rapidly and appear to jump. Venus at maximum brightness, low on the horizon,
through turbulent air, is the single most reported UFO in history and has been chased by
military aircraft on multiple continents.

\textbf{Meteors and re-entries} deserve their own note. A bolide can be brighter than the full
moon, fragment visibly, appear to travel slowly and horizontally when seen end-on, and be
reported as a formation of lit windows in a structured craft --- the 1968 Zond IV re-entry over the
eastern United States produced exactly that report from multiple witnesses, described as a
cigar-shaped craft with square windows. It was a rocket body breaking up.

\begin{funfact}
And then there are the ones that are not weather at all. Starlink satellite trains produce a
perfect line of moving lights and generated a global spike in reports from 2019. Chinese lanterns
produce orange orbs in silent formation that drift with the wind and wink out --- the 2007
Alderney and countless UK ``fleet'' sightings are lanterns, and the correlation with wedding
season is not subtle. High-altitude research balloons at 100,\!000 feet catch sunlight for an hour
after dark. Google's Project Loon put balloons in the stratosphere for years. If a thing in the
sky is silent, slow, and drifting downwind, the wind is the identification.
\end{funfact}

\subsubsection*{Falls, jellies and lights: the rest of the catalogue}

Beyond craft and beyond weather, the sky produces a catalogue of genuinely odd reports, and I want
to run through them because the pattern across the category is instructive.

\textbf{Falls of fish and frogs} are real and documented since antiquity, with modern cases
verified by authorities --- Honduras has an annual event, and there are well-attested cases from
England, Japan and Australia. The mechanism is waterspouts and tornadic updraught lifting the
contents of a shallow pond and depositing them downwind, which explains why the animals are
typically all one species and one size class: the pond sorted them.

\textbf{Coloured rain} --- red rain in Kerala in 2001 --- was analysed and found to contain spores
of a local algae, \emph{Trentepohlia}. It had been proposed as extraterrestrial by Chandra
Wickramasinghe, who has proposed that about a great many things.

\textbf{Angel hair} is spider silk, from mass ballooning dispersal events, which coat landscapes
in filament and drift down from great height.

\textbf{Star jelly} is variously slime mould, cyanobacteria, or the oviducts of frogs deposited by
predators, all of which have been identified in specific cases.

\textbf{Ice falls} --- megacryometeors --- are real, up to several kilograms, and mostly but not
entirely attributable to aircraft; some fall from clear skies and remain argued over.

\textbf{Airborne debris} accounts for much more than people expect: pieces of aircraft, cargo,
balloon payloads, and space junk. Around 100 tonnes of meteoritic material arrives daily, mostly
as dust.

\textbf{The Hessdalen lights} in Norway deserve better than a paragraph. Recurring luminous
phenomena in a specific valley, under continuous instrumented observation since 1984 with
cameras, spectrometers, radar and magnetometers --- an ongoing scientific field study of an
anomalous light phenomenon. Results suggest a plasma or combustion process possibly involving
valley geology and dust, with some events correlating with radar returns. It is unresolved.
\emph{This} is how you study an anomaly: pick a place where it recurs, put instruments there, and
wait for years. Compare it to the field's usual method of interviewing people about something
that happened once, decades ago.

\textbf{Earthquake lights} were dismissed for centuries and are now accepted, with a proposed
mechanism involving stress-induced charge separation in certain rock types. Another case where the
witnesses were right.

\textbf{Tunguska} is the one that gets folded into ufology, and it belongs here because of how it
is misused rather than because of what it was. On the morning of 30 June 1908 something detonated
several kilometres above the Podkamennaya Tunguska river in Siberia. The airburst flattened
roughly 770 square miles of forest and some eighty million trees, registered around 5.0 on
seismographs, and sent atmospheric pressure waves that were recorded on barometers in London
thousands of kilometres away. Estimates of the yield run in the range of ten to fifteen megatons,
which is to say something on the order of a thousand times the Hiroshima device. Crucially, there
is no impact crater and no large recovered meteorite, because the body broke up under
aerodynamic stress before reaching the ground.

That absence of a crater is the entire foothold for the exotic readings --- a crashing alien craft,
a natural nuclear explosion, an antimatter fragment, a small black hole. None of them is needed.
An airburst is the expected behaviour of a stony asteroid or comet fragment of that size entering
at that speed, and it is exactly what happened on a smaller scale over Chelyabinsk on 15 February
2013, an event captured by hundreds of dashboard cameras and analysed to death, which produced
the same signature: bright bolide, high-altitude detonation, blast damage over a wide area, no
crater, and only small recovered fragments. Tunguska is not a UFO case. It is the best available
demonstration that the sky is genuinely dangerous without any help from anybody, and the reason I
include it is that the field's instinct on encountering an unexplained gap --- reach for the most
exotic filler --- is precisely the instinct this course exists to break.

\begin{srcnotes}
\item McDowell, J.\ (10 December 2009). Commentary on the Norwegian spiral and the Bulava third-stage
failure, \emph{Jonathan's Space Report} and contemporaneous press interviews.
\item Russian Ministry of Defence (10 December 2009). Statement confirming failed Bulava test launch
from the White Sea.
\item National Astronomical Observatory of Japan / Subaru--Asahi Star Camera (18 January 2023).
Time-lapse of the Falcon~9 fuel-dump spiral over Maunakea.
\item Geophysical Institute, University of Alaska Fairbanks (15 April 2023). Explanation of the
spiral observed with the aurora over interior Alaska.
\item Gindilis, L.\ M.\ et al.\ (1979). Observations of Anomalous Atmospheric Phenomena in the
USSR, Academy of Sciences of the USSR --- the Petrozavodsk investigation.
\end{srcnotes}

% ---------------------------------------------------------------
\section{Dyatlov Pass}\label{sec:dyatlovpass}

Every thread so far now converges on one incident, which is why it sits here rather than in a
chapter of its own. Dyatlov Pass contains lights in the sky
(Section~\ref{sec:paranormalskies}), missing soft tissue on a body with no external wound
(Section~\ref{sec:cattlemutilations}), a radiological detail that sounds like a weapons programme
(Section~\ref{sec:philadelphiaexperiment}), and an official finding written in language
deliberately chosen to say nothing. It is the first case in the chapter where the individual
explanations are not merely correct one at a time but have to be assembled, in order, to account
for a single night. Do that assembly carefully here, because it is the exact exercise
Section~\ref{sec:skinwalker} will ask you to perform without assistance.

In early February 1959, nine experienced ski hikers from the Ural Polytechnic Institute, led by
Igor Dyatlov, died on the eastern slope of Kholat Syakhl in the northern Urals. The details are
why the case has never gone away.

\figwrap[r]{0.30}{The tent on the eastern slope of Kholat Syakhl as the search party found it on
26 February 1959, photographed for the case file. It had been cut open from the inside, and most of
the party left it in socks or barefoot. This is the whole of the physical scene the legend rests on.}{}{fig:dyatlovtent}%
Their tent was found cut open from the inside. All nine had left it in severe cold ---
something like $-25$ to $-30^\circ$C --- inadequately dressed, several barefoot or in socks. Five
died of hypothermia. Four, found later in a ravine under deep snow, had severe internal injuries:
Thibeaux-Brignolle had a major skull fracture, Dubinina and Zolotaryov had multiple crushed ribs,
with no corresponding external wounds --- the investigator's phrase was that the force resembled
``a car crash.'' Dubinina was missing her tongue and eyes. Some clothing showed elevated
radioactivity. Some bodies had an orange-brown discolouration. Other hikers in the region
reported orange spheres in the sky that month. The Soviet inquiry closed with the formulation
that a ``compelling natural force'' had caused the deaths, and the area was closed for three
years.

The explanations, in ascending order of plausibility:

\textbf{Avalanche} was long dismissed because the slope angle seemed too shallow and there was no
obvious debris. In 2019 the Russian authorities reopened and closed the case with a slab avalanche
finding, and in 2021 Puzrin and Gaume published a modelling paper in \emph{Communications Earth
\& Environment} showing that cutting into the slope to pitch the tent, combined with persistent
katy winds loading the slab, could produce a delayed, small, localised slab release capable of
causing exactly those crush injuries in people lying in a tent. The injuries match the biomechanics
of compression under a slab against a hard surface, which also explains the absence of external
wounds.

\textbf{Paradoxical undressing} and \textbf{terminal burrowing} are established hypothermia
behaviours: as core temperature falls, peripheral vasodilation produces a false sensation of heat
and victims remove clothing, then seek enclosed spaces. This accounts for the state of the
bodies and for the ravine.

\textbf{The soft-tissue loss} is scavenging and decomposition in running water --- the ravine
bodies were found in a stream bed after months, and eyes, tongue and lips are the first tissues
lost, as Section~\ref{sec:cattlemutilations} explained.

This is the point of putting Dyatlov Pass in this chapter at all, so it is worth setting the two
lists side by side (Figure~\ref{fig:dyatlovmutil}). The tissue reported missing from Dubinina and
Zolotaryov is not a random selection. It is the top of the same ranked list that the cattle cases
produce, in the same order, for the same reason: eyes, tongue and the margins of the mouth are the
wettest, softest and least protected tissue on a body, they are exposed first, and they are what
insects, small scavengers and running water take first. What makes the case rare is not the
mechanism. It is that this is one of the very few times the mechanism has acted on human remains
with a criminal file open on them, an autopsy performed, and a state's word for it that the cause
could not be named. Given a cow in a field, the same findings are read as post-mortem scavenging.
Given nine young people and a closed inquiry, they are read as mutilation.

\Needspace*{3.3in}
\begin{center}
\refstepcounter{figph}\label{fig:dyatlovmutil}%
\begin{figpanel}[width=0.96\textwidth]
\centering
{\setlength{\fboxsep}{0pt}\setlength{\fboxrule}{0.5pt}%
\fcolorbox{ufogreen}{white}{\adjustimage{max size={0.94\textwidth}{3.9in}}{dg_dyatlovmutil.png}}}\par\vspace{3pt}
{\RaggedRight\scriptsize\color{steel}\textbf{\textcolor{ufogreenink}{Figure \thefigph.}}
What is missing, and in what order it goes. The tissue absent from the ravine bodies at Dyatlov
Pass is the top of the same list as the cattle mutilation reports in
Section~\ref{sec:cattlemutilations}, which is what the autopsy findings and the scavenging
literature both predict for remains left in a snow-filled stream bed for three months. The
autopsy photographs of Dubinina and Zolotaryov exist in the case file and are deliberately not
reproduced here: they are graphic images of identifiable dead people whose families are still
living, and reprinting them would add nothing the findings do not already say.\par
\vspace{1pt}{\color{midgrey} Author's own diagram. Compiled from the 1959 autopsy findings and the
scavenging literature cited in Section~\ref{sec:cattlemutilations}.}\par}
\end{figpanel}
\end{center}
\vspace{6pt}

\textbf{The radioactivity} was on the clothing of members who had worked with radioactive
materials at the institute, which the file records.

\textbf{The orange spheres} are almost certainly R-7 rocket launches from Baikonur or Plesetsk ---
the timing and geography fit, and rocket plumes in the upper atmosphere produce exactly that.

\begin{funfact}
I find Dyatlov Pass the most sobering case in this book, and not because of the mystery. Nine
young people died in the cold, badly, and their families spent sixty years being told the file
was closed by a state that closed files as a matter of habit. The secrecy produced the legend ---
not because there was a cover-up of aliens, but because a government that will not explain
anything teaches its citizens that nothing is explained. Chapter~\ref{ch:government} is about the American version
of the same lesson.
\end{funfact}

% ---------------------------------------------------------------
\section{Cryptozoology}\label{sec:cryptozoology}

Dyatlov Pass had bodies, a tent, an autopsy and a file --- too much evidence, badly handled.
This section is the opposite failure, and I have put it immediately afterwards on purpose.
Here the evidence never arrives at all, decade after decade, and the field carries on regardless.
It is also the one comparison in the book that is genuinely fair to ufology, because
cryptozoology is looking for something that would be made of ordinary matter, would live in an
ordinary habitat, and would leave ordinary remains --- and it still cannot produce one.

Cryptozoology is the search for animals whose existence is unconfirmed, and it belongs in this
book for one reason: it is the closest available control group. It uses the same evidence types as
ufology, faces the same problems, and --- crucially --- it occasionally \emph{wins}, which lets us
see what a win looks like.

The wins are real and worth knowing. The okapi was confirmed in 1901, the mountain gorilla in
1902, the Komodo dragon in 1912, the coelacanth in 1938 after being presumed extinct for 66
million years, the saola in 1992, and the giant squid was photographed alive only in 2004 and
filmed in 2012. New species are described continuously --- thousands a year --- including large
vertebrates. Anybody who tells you biology is finished is wrong.

Notice, though, what confirmation required in every case: a body, a specimen, a type description,
and a place in the taxonomic literature. Not footage. Not testimony. A carcass.

Which is what the famous cryptids lack.

\figwrap[r]{0.38}{Frame 352 of the Patterson--Gimlin film: the most analysed and least conclusive footage in cryptozoology.}{https://commons.wikimedia.org/wiki/File:Patterson-Gimlin_film_frame_352.jpg}{fig:pattersongimlin}
\textbf{Bigfoot.} The Patterson--Gimlin film of 1967 (see Figure~\ref{fig:pattersongimlin}) remains the centrepiece, and after nearly
sixty years of analysis it is unresolved in the sense that neither side can force the other's
hand. Bob Heironimus claimed to have worn the suit. What is not unresolved: no body, no bones, no
teeth, no scat with novel DNA, no roadkill, and --- this is the killer --- no fossil record of
any great ape in the Americas, ever. A breeding population of large primates requires thousands of
individuals over thousands of years, and thousands of individuals die. A 2014 study by Sykes and
colleagues sequenced 30 alleged samples: every one was a known animal, mostly bear, and two
supposed Yeti samples were Himalayan brown bear.

\textbf{Loch Ness.} Sonar surveys, including comprehensive sweeps, found nothing large. A 2018
environmental DNA survey sampled the entire loch and found no reptilian DNA at all --- and a great
deal of eel DNA, which produced the current best hypothesis and a rather charming one. The 1934
``surgeon's photograph'' was admitted in 1994 to be a toy submarine with a model head. And the
loch is oligotrophic: its food web cannot support a large predator population.

\textbf{Thylacine}, by contrast, is the interesting one. It genuinely existed, the last confirmed
individual died in Hobart in 1936, and sighting reports continue. That is a case where the animal
is real, extinction is recent, habitat exists, and the search is legitimate --- and it still has
not produced a specimen in ninety years, which tells you how hard producing a specimen is even
when the animal was definitely there.

\begin{funfact}
The lesson cryptozoology teaches ufology is about \emph{trail cameras}. There are now millions of
motion-triggered cameras in North American forests, operating continuously for two decades, and
they have produced an extraordinary photographic record of wildlife --- including rare, shy,
nocturnal animals in enormous numbers. Not one unambiguous Sasquatch. Meanwhile the same interval
put a high-resolution camera in every pocket on Earth, and the quality of UFO photographs has
not improved. When the observing capacity increases by orders of magnitude and the evidence does
not, that is data.
\end{funfact}

\subsubsection*{Marble Hornets, and a cryptid with a birth certificate}

Everything in this section suffers from the same problem: the origin of the legend is lost. Nobody
can interview the first person to see a Sasquatch, because the first person to see one is
unreachable in every case --- either they are dead, or the story is older than the paperwork, or
the trail runs into oral tradition and stops. That is precisely why the next case matters. There
is exactly one modern cryptid whose origin is documented to the hour, and it lets us watch the
whole process run in daylight.

On 10 June 2009, a user named Victor Surge --- Eric Knudsen --- posted two black-and-white
photographs to a Photoshop thread on the Something Awful forums. Children in a playground, and
behind them a very tall thin figure in a dark suit with no face. He added a few lines of invented
caption in the voice of an archive. Other users immediately began adding their own photographs,
their own fragments of witness testimony, their own dates. Within weeks the figure had a name, a
behaviour --- it takes children, it appears in the trees, it cannot be photographed cleanly ---
and a set of rules nobody had written down as rules.

What turned a forum post into a folklore organism was \emph{Marble Hornets}. Beginning in June
2009, Troy Wagner and Joseph DeLage released a YouTube series presented as recovered footage: a
student finds tapes from an abandoned film project and uploads them one at a time, and the
something in the tapes is never explained. It ran to eighty-seven entries and a parallel
adversarial channel, \emph{totheark}, posting cryptic replies, and it finished in 2014. It is
worth watching for its method rather than its story, because the method is the entire subject of
this book. The footage is deliberately degraded. The figure is never in focus and never centred.
The tapes arrive out of order and incomplete. The narrator's own reliability decays. Nothing is
asserted; everything is implied by a gap in the record. Those are the exact properties of the
footage that ufology treats as its strongest evidence --- and here they were engineered on
purpose, by two people, on a consumer camcorder, for nothing.

The consequences were not confined to the internet. The legend acquired ritual: children's games,
summoning rules, a canon of sightings. In May 2014 two twelve-year-old girls in Waukesha,
Wisconsin, stabbed a classmate nineteen times and told police they had done it for Slenderman. She
survived. It is the clearest demonstration in this book that a thing does not have to exist in
order to produce a body.

\begin{funfact}
Slenderman is the only entry in the roster that follows with a known author, a known date and a
known motive --- and it is also the only one whose folklore we can prove developed exactly the way
folklorists say folklore develops: collaboratively, anonymously, and with each contributor
sincerely believing they were reporting rather than inventing. Ufology has no case with this
provenance. That is not a compliment to ufology. It means we have a working model of how a
sighting-based mythology assembles itself out of nothing, and the model fits the older cryptids
uncomfortably well.
\end{funfact}

% ===================== CRYPTID TAXONOMY =====================
% The control-group roster: sixty animals that are claimed to exist and have
% never been produced. Six cards to a page, three columns by two rows, over
% the night-swamp plate rather than the starfield, so this list and the Alien
% Taxonomy in Chapter 6 can be read against each other.
% Entries are Evan's own list, sorted A--Z inside each of his seven groups,
% with the groups in his order. The group name opens the italic line of every
% card so the taxonomy stays legible without breaking the six-card grid.
% Portraits for the newer entries are placeholders and will be replaced plate
% for plate.

So here is the control group, on the same card system as the alien roster and
for the same reason. Read the two lists side by side and notice how similar the
evidence is: a film nobody can resolve, a photograph later admitted, a
footprint cast by an enthusiast, a witness with no reason to lie. Then notice
the thing this chapter is actually about --- how many of these animals stand
up. Bears do not walk to the horizon on two legs. Apes do not either. Yet the
reports insist on it, over and over, on every continent, in cultures with no
contact between them. That is not a fact about biology. It is a fact about us,
and it is the same fact that shapes the alien roster: we describe the unknown
by standing it upright and giving it a face.

The final line of each card is the part that matters. In almost every entry
there is a perfectly ordinary animal standing behind the report, and the few
where there is not are the few worth a researcher's evening.

\swampfieldon
\rosterwidesix

\begin{multicols}{3}[\rosterhead{Cryptid Taxonomy}]

% ---------- 1. HOMINIDS & PRIMATES (YETI-LIKE) ----------

\cryptidcard{cryp_agogwe}{Agogwe}{Hominid; East Africa, reported
1900--1930s}{A small woolly bipedal humanoid, three to five feet tall, reported
from Tanzania and Mozambique by colonial hunters. The accounts are brief and
second-hand, and they arrive in exactly the decades when European hunters were
first writing up East African forest. Small upright figures at the edge of
vision are the oldest error in the file.}{A juvenile chimpanzee or a bushbaby
seen at a distance, or a spotted hyena's cub standing.}

\cryptidcard{cryp_almas}{Almas}{Hominid; Caucasus and Pamir, reported for
centuries}{A wild man --- hairy, silent, upright --- filling the same slot in
Central Asian folklore that the Yeti fills further south. The interest here is
archaeological rather than zoological: the region did hold other hominins
within the memory of our species, and folklore can carry a real observation for
a very long time. It cannot carry a specimen.}{A bear on its hind legs, or an
ancestral memory of something gone for forty thousand years.}

\cryptidcard{cryp_batutut}{Batutut}{Hominid; Vietnam and Laos, reported
1960s--1970s}{The \emph{Nguoi Rung}, or forest man: a hairy bipedal hominid
reported by soldiers and by Vietnamese zoologists in the Annamite jungle. The
range is genuinely under-surveyed and has yielded new large mammals in living
memory, which is the strongest argument any hominid entry has. It has still
yielded no hominid.}{A sun bear or a gibbon walking upright, both of which do
so readily on the ground.}

\cryptidcard{cryp_bigfoot}{Bigfoot (Sasquatch)}{Hominid; Pacific Northwest,
named 1958}{The type specimen of modern cryptozoology: a large ape-like
hominid in temperate North American forest. A breeding population needs
thousands of animals, and thousands of animals leave bodies, bones, scat and
roadkill. Two decades of continuous trail cameras across the range have
produced none of it.}{Black bears, misjudged scale, and a track record of
admitted hoaxes including the original casts.}

\cryptidcard{cryp_mapinguari}{Mapinguari}{Hominid; western Amazon, reported
from the nineteenth century}{A hulking, foul-smelling forest giant of Amazonian
tradition, sometimes read by researchers as folk memory of a ground sloth. The
sloth reading is attractive and unprovable: \emph{Megatherium} was real, it was
here, and it went out roughly when people arrived. Memory is not evidence of
survival.}{A giant anteater rearing to threat-display, which is upright, tall,
clawed and rank.}

\cryptidcard{cryp_maricoxi}{Maricoxi}{Hominid; Mato Grosso, Brazil, reported
1914}{Aggressive ape-men, bow-armed and grunting, described by Percy Fawcett in
the Amazon. This one is a lesson in provenance rather than in zoology: the
source is a single celebrated explorer writing for an audience, in a career
that ended in a disappearance the public wanted to be a mystery.}{An
unacquainted Indigenous group, described in the language of an era that
routinely dehumanised them.}

\cryptidcard{cryp_orangpendek}{Orang Pendek}{Hominid; Sumatra, reported since
the 1910s}{The short man: a ground-dwelling bipedal primate of the Kerinci
Seblat rainforest, reported by loggers, villagers and a number of trained
field biologists. It is the best of the hominid cases --- plausible size,
plausible habitat, credible witnesses --- and after a century of camera traps
and hair samples it still has no type specimen.}{A sun bear, an agile gibbon,
or an undescribed orangutan population.}

\cryptidcard{cryp_yeren}{Yeren}{Hominid; Hubei and Shennongjia, China,
reported since the 1940s}{A reddish-furred wild man of the Chinese mountains,
supported by state-backed expeditions in the 1970s and 1980s. The hair samples
that came back were the usual ones. The expeditions themselves are the
interesting artefact: this is what a well-funded, officially sanctioned search
for a hominid produces, which is nothing.}{Golden snub-nosed monkeys, bears,
and hair later identified as human or bovine.}

\cryptidcard{cryp_yeti}{Yeti (Abominable Snowman)}{Hominid; Himalaya, in the
West since 1921}{The white-furred high-altitude hominid, and the case with the
best physical trail --- relics, scalps, hides and hair held in monasteries and
museums. That trail is why it belongs here: nine separate samples have been
sequenced, and the genetics came back Asian black bear, brown bear, Tibetan
brown bear and dog.}{Himalayan brown and black bears; tracks widened by melt
and refreeze.}

\cryptidcard{cryp_yowie}{Yowie}{Hominid; eastern Australia, colonial reports
from the nineteenth century}{The Australian hairy man, reported down the Great
Dividing Range, and the entry with the hardest biogeography on the list.
Australia has never had a native primate of any kind in any period, so a
breeding population of large apes needs an arrival, a history and a fossil
record. There is none.}{Feral animals, kangaroos standing upright, and hoaxes
with a long local pedigree.}

% ---------- 2. LACUSTRINE & MARINE MONSTERS ----------

\cryptidcard{cryp_bessie}{Bessie (South Bay Bessie)}{Lake monster; Lake
Erie, USA, reported since 1793}{A serpentine animal in the eastern basin of
Lake Erie, reported from shore and from boats for two centuries. Erie is the
shallowest of the Great Lakes and the warmest, which makes it the one most
likely to hold a large sturgeon and the one least likely to hide anything for
two hundred years.}{Lake sturgeon, which reach seven feet, and standing waves
in a lake famous for short steep chop.}

\cryptidcard{cryp_brosno}{Brosno Dragon}{Lake monster; Lake Brosno, Russia,
folklore and 1940s claims}{A reptilian animal in a lake west of Moscow, said in
the best-known version to have swallowed a German aircraft during the Second
World War. The story is a good example of a folklore creature acquiring a
modern, dateable, checkable event --- and of the event turning out to be
unsourced when checked.}{Gas eruptions from the lake bed, which surveys of
Brosno have in fact recorded, throwing up mounds of water.}

\cryptidcard{cryp_cadborosaurus}{Caddy (Cadborosaurus)}{Sea serpent; Pacific
coast of North America, named 1933}{A long horse-headed sea serpent reported
from Alaska to Oregon, with a supposed carcass photographed at Naden Harbour in
1937 and then lost. The lost specimen is the whole case. A century of coastal
reports and the only physical evidence went into a rendering plant without
anyone measuring it.}{Basking sharks in decay, oarfish, and lines of sea lions
surfacing in file.}

\cryptidcard{cryp_champ}{Champ}{Lake monster; Lake Champlain, reports since
1819}{The North American Nessie, with a formal reward, a state resolution
protecting it, and the 1977 Mansi photograph --- the best lake-monster picture
ever taken, and one whose provenance cannot be checked because the negative and
the location were never produced.}{Lake sturgeon, floating logs, and standing
waves in a long narrow lake.}

\cryptidcard{cryp_kraken}{Kraken}{Sea monster; North Atlantic, Norse accounts
from the twelfth century}{The one that came true, and therefore the most
important card on this page. A ship-sinking tentacled monster of Scandinavian
sailors' lore turned out to be a real animal: \emph{Architeuthis}, the giant
squid, forty feet long, finally filmed alive in 2012. What made it science was
a body, then many bodies.}{The giant squid --- and this entry is filed as
solved, not open.}

\cryptidcard{cryp_lochness}{Loch Ness Monster (Nessie)}{Lake monster; Loch
Ness, Scotland, modern era from 1933}{The most investigated cryptid on Earth:
sonar sweeps, submersibles, and in 2018 an environmental-DNA survey that
sampled the whole water column. The eDNA found eels, plenty of them, and no
reptile. The 1934 ``surgeon's photograph'' was confessed as a toy submarine
with a putty neck.}{European eels, boat wakes, and deer or otters swimming.}

\cryptidcard{cryp_lusca}{Lusca}{Marine; Bahamian blue holes, reported since
the nineteenth century}{A giant predatory cephalopod said to inhabit the blue
holes of the Bahamas and to pull swimmers down. The blue holes are real, deep,
tidal and genuinely dangerous: they breathe in and out with the tide hard
enough to drown a diver, which is a sufficient monster.}{Tidal suction in
karst holes, plus large Caribbean reef octopus and moray eels.}

\cryptidcard{cryp_memphre}{Memphre}{Lake monster; Lake Memphremagog,
Quebec--Vermont, reports since 1816}{A reptilian animal in a lake straddling
the Canada--US border, with a two-century reporting record kept unusually well
by local historians. The archive is the value here: read in order, the sightings
follow the tourist season and the newspaper coverage more closely than they
follow the water.}{Sturgeon, otters in line, and wave interference between two
steep shores.}

\cryptidcard{cryp_mokelembembe}{Mokele-Mbembe}{Sauropod claim; Congo River
basin, Western reports from 1776}{A long-necked reptile in the swamps of the
Likouala, sought by more than fifty expeditions, several explicitly funded to
prove that dinosaurs survived. Sauropods were not swamp animals, the Congo
basin has a Cenozoic fossil record with nothing in it, and no expedition has
returned with a photograph.}{Elephants swimming with the trunk raised, and
rhinoceros folklore carried inland by trade.}

\cryptidcard{cryp_ningen}{Ningen}{Marine humanoid; Antarctic waters, internet
reports from the 2000s}{A vast pale humanoid said to have been seen from
Japanese research vessels near the ice edge. This card is here as a control on
the others: it has no witness, no date, no vessel and no log entry, and it
begins on a message board in 2007. It shows how fast a cryptid can now be
manufactured.}{Ice floes and internet folklore; there is no case to explain.}

\cryptidcard{cryp_ogopogo}{Ogopogo}{Lake monster; Okanagan Lake, British
Columbia, Syilx accounts and settler reports from 1872}{A serpentine animal in
a deep glacial lake, known to the Syilx people as \emph{n'ha-a-itk} long before
settlers named it after a music-hall song. This one is personal: my mother grew
up in the Okanagan, and I have stood at the lakeside monument in Kelowna many
times. Standing there is exactly how you learn what the lake does --- how a
boat wake crossing a flat evening surface throws three black humps that hold
their shape for half a minute.}{Wake interference on a long deep lake, plus
sturgeon, logs and otters in file.}

\cryptidcard{cryp_seaape}{Steller's Sea Ape}{Marine mammal; Shumagin
Islands, one sighting, 1741}{A playful unclassified mammal watched for two
hours by Georg Steller from the \emph{St.\ Peter}. Steller is the best witness
in this book --- a trained naturalist who described the sea cow, sea lion and
eider that bear his name, all of them confirmed. One animal in his notebook
never was, and that is what a single unrepeated observation is worth.}{A
northern fur seal in its first year, or a juvenile sea otter at play.}

% ---------- 3. WINGED CRYPTIDS (AVIAN & CHIROPTERAN) ----------

\cryptidcard{cryp_ahool}{Ahool}{Winged; Java, Indonesia, reported from the
1920s}{A bat the size of a man, named for its call, said to hunt fish along
jungle rivers at night. Every element belongs to an animal that is already there
and already remarkable, which is usually the shape of these cases: the witness
has seen something real and estimated its size in the dark, over water, with
nothing to scale it against.}{The Javan flying fox, whose wingspan reaches five
feet and whose face is startling at close range.}

\cryptidcard{cryp_jerseydevil}{Jersey Devil}{Winged; Pine Barrens, New Jersey,
legend from 1735, wave in 1909}{A kangaroo-bodied, bat-winged, horse-headed
thing with a colonial curse story attached. The January 1909 wave is the useful
part: a week of sightings across two states, hoof-prints in snow, schools
closed --- and at the end of it a Philadelphia showman admitting he had glued
bronze wings to a kangaroo.}{Sandhill cranes, owls, and a documented
promotional hoax.}

\cryptidcard{cryp_kongamato}{Kongamato}{Winged; Jiundu swamps, Zambia,
reported from the 1920s}{A tooth-billed flying reptile of the Kaonde, whose name
means breaker of boats. Frank Melland's 1923 book is the origin of the
pterosaur reading, and it arrived by a route worth noticing: he showed
villagers a picture of a pterosaur and asked whether that was it.}{Saddle-billed
storks and large fruit bats, plus a leading question that supplied its own
answer.}

\cryptidcard{cryp_mothman}{Mothman}{Winged humanoid; Point Pleasant, West
Virginia, 1966--67}{A grey winged humanoid with red eyes, seen for thirteen
months and then bound forever to the collapse of the Silver Bridge in December
1967. The bridge failed because of a single corroded eyebar. The prophecy
reading was built afterwards, from newspaper files, by a writer who said so.}{A
sandhill crane far off migration, whose red eye-patch and seven-foot span fit
the first reports exactly.}

\cryptidcard{cryp_owlman}{Owlman}{Winged humanoid; Mawnan, Cornwall, England,
1976}{A large grey owl-like figure seen above a church tower by two girls on
holiday. The investigator who collected the case, Tony ``Doc'' Shiels, was a
professional magician who also produced one of the famous Nessie photographs of
the same period, which is not proof of a hoax but is the first thing an honest
file records.}{A eurasian eagle-owl, five feet across the wings, on a
silhouetted tower at dusk.}

\cryptidcard{cryp_ropen}{Ropen}{Winged; Umboi Island, Papua New Guinea,
expeditions from the 1990s}{A nocturnal flying reptile said to glow in flight,
sought by expeditions funded largely to demonstrate that pterosaurs survived.
The bioluminescence claim is the tell: no vertebrate produces sustained
bioluminescence in air, and the local accounts describe a light, not an
animal.}{Frigatebirds and flying foxes, plus bioluminescent fungus and
misidentified meteors.}

\cryptidcard{cryp_snallygaster}{Snallygaster}{Winged; Frederick County, Maryland,
German-settler lore, wave in 1909}{A bird-reptile with a metal beak and
tentacles, from \emph{Schneller Geist}, quick spirit. Its 1909 revival was
printed in the \emph{Middletown Valley Register} by editors who later admitted
they were inventing instalments to sell papers, and the Smithsonian was said to
have offered a reward.}{A newspaper circulation stunt built on genuine
eighteenth-century immigrant folklore.}

\cryptidcard{cryp_thunderbird}{Thunderbird}{Winged; across North America,
Indigenous tradition and settler reports}{An enormous bird whose wingbeats make
thunder, held sacred across many nations and then flattened by settlers into a
monster-bird story with a fake photograph attached --- the ``Tombstone
thunderbird'' picture that everyone remembers seeing and nobody can produce.}{
California condors and white pelicans, both with wingspans near ten feet, plus
a false memory of a photograph that does not exist.}

\cryptidcard{cryp_vanmeter}{Van Meter Monster}{Winged; Van Meter, Iowa, one
week in 1903}{A horned winged figure that townspeople said projected a blinding
beam from its forehead and smelled foul. Five respected men fired at it. The
case is unusually well bounded --- one town, one week, named witnesses, local
press --- and it stops as abruptly as it starts.}{Bats disturbed from an
abandoned coal mine, seen against lantern glare during a regional bat-story
season.}

% ---------- 4. TERRESTRIAL MAMMALS & QUADRUPEDAL PREDATORS ----------

\cryptidcard{cryp_beastgevaudan}{Beast of G\'evaudan}{Predator; Margeride,
France, 1764--1767}{The best-documented cryptid in history, and the one that
killed: roughly a hundred people dead in three years, church records, royal
huntsmen, a national panic. It ends the way a real case ends --- an animal was
shot, the deaths stopped, and the carcass was examined and paraded.}{One or
more very large wolves, possibly wolf-dog hybrids, hunting people during a
famine.}

\cryptidcard{cryp_chupacabra}{Chupacabra}{Predator; Puerto Rico, 1995, then
the Americas}{The goat-sucker: first described in Puerto Rico as a spined
bipedal thing with red eyes, later reported across the southwest as a hairless
dog. The Puerto Rican original was traced to a witness describing the creature
from the film \emph{Species}, released weeks earlier, and every carcass
submitted since has been sequenced.}{Coyotes and dogs with advanced sarcoptic
mange, which strips the coat and swells the skin.}

\cryptidcard{cryp_dingonek}{Dingonek}{Predator; rivers of East Africa,
reported 1907}{A scale-armoured, walrus-tusked river predator described by the
hunter John Alfred Jordan in the Maggori River, Kenya. It is the standard
colonial-hunter cryptid: a single armed European, a river he did not know, one
glimpse, and a printed account with no specimen and no second witness.}{A
hippopotamus seen head-on in turbid water, or a large Nile crocodile.}

\cryptidcard{cryp_easternpuma}{Eastern Puma}{Predator; eastern North America,
reported continuously since being declared extinct}{The ghost cat: cougars
persisting east of the Mississippi after the subspecies was formally declared
extinct in 2011. This is the one entry on the roster where the sceptics were
partly wrong --- confirmed animals do turn up in the east, and DNA has traced
them to dispersing western males and escaped captives.}{Confirmed dispersing
western cougars, released pets, and bobcats or large domestic cats misjudged
for size.}

\cryptidcard{cryp_gef}{Gef the Talking Mongoose}{Anomalous animal;
Cashen's Gap, Isle of Man, 1931--1939}{A talking mongoose that lived in the
walls of the Irving farmhouse, gave interviews, and was investigated in person
by Harry Price. It matters because it was taken seriously by the press of a
literate country for eight years, on the testimony of one family, with hair
samples that came back as the family dog.}{A family hoax, most probably driven
by the daughter, sustained by national newspaper attention.}

\cryptidcard{cryp_nandibear}{Nandi Bear}{Predator; western Kenya, reported
1900--1930s}{A ferocious nocturnal carnivore of the Nandi escarpment, described
as high in the shoulder and low at the hip. That silhouette is the whole entry:
it is the outline of an animal that certainly lives there, and the reports drop
away as the region is surveyed.}{The spotted hyena, whose sloping back, nocturnal
habits and bone-cracking bite fit every account.}

\cryptidcard{cryp_ozarkhowler}{Ozark Howler}{Predator; Ozark Mountains, USA,
folklore and modern reports}{A bear-sized black cat-canine with horns and a cry
like a woman screaming. The modern history is documented rather than
speculative: a folklorist admitted seeding Ozark Howler stories on the early
internet to see how far a manufactured cryptid would travel, and it
travelled.}{A mountain lion's scream, a black bear at night, and a documented
online hoax layered onto older hill folklore.}

\cryptidcard{cryp_queenslandtiger}{Queensland Tiger}{Predator; north
Queensland, Australia, reported into the twentieth century}{A striped cat-like
marsupial predator, the \emph{yarri} of Aboriginal accounts, reported by
graziers and by at least one naturalist. Australia did hold a marsupial lion,
\emph{Thylacoleo}, and it did have stripes in some reconstructions --- and it
left a fossil record that stops.}{Feral cats grown large, quolls, and the
surviving memory of \emph{Thylacoleo carnifex}.}

\cryptidcard{cryp_shunkawarakin}{Shunka Warakin}{Predator; Montana, USA, shot
1886}{A hyena-like animal that killed livestock, shot near Sun River and mounted
by a taxidermist --- and then, unusually, the mount survived. It was rediscovered
in 2007, sampled, and sequenced. This is what closure looks like when a body
exists.}{A wolf, confirmed by DNA from the surviving mount, with an unusual
sloped-back conformation.}

% ---------- 5. SERPENTINE & DRACONIC REPTILES ----------

\cryptidcard{cryp_altamahaha}{Altamaha-ha}{River monster; Altamaha River,
Georgia, USA, reports from the eighteenth century}{A sturgeon-headed,
snake-bodied animal in the tidal reaches of the Altamaha, with Tama Indian
accounts and steady modern sightings near Darien. The candidate here is not a
guess: Atlantic sturgeon spawn in this exact river and reach fourteen
feet.}{Atlantic sturgeon, whose armoured back and rolling surface leaps match
the descriptions almost line for line.}

\cryptidcard{cryp_bunyip}{Bunyip}{Amphibious; swamps and billabongs, Australia,
Aboriginal tradition}{A marsh-dwelling animal with a dog-like face, deep in the
traditions of many Aboriginal nations and recorded by settlers from the 1810s.
It is also the one case where colonists brought bones to Aboriginal people for
identification and were told, correctly, which animal they belonged to.}{The
Australian fur seal or leopard seal moving up rivers, plus surviving memory of
\emph{Diprotodon} and other extinct megafauna.}

\cryptidcard{cryp_buru}{Buru}{Reptile; Rilo Valley, Assam, India, accounts
recorded 1945--47}{A large aquatic reptile said by the Apatani to have lived in
the marshes of their valley until their ancestors drained them for rice.
Charles Stonor and Ralph Izzard investigated on the ground and found a coherent
oral history --- and a valley that had genuinely been drained, so that whatever
the animal was, its habitat was gone before science arrived.}{A large monitor
lizard or crocodilian, extirpated locally by wetland drainage.}

\cryptidcard{cryp_hoopsnake}{Hoop Snake}{Folklore serpent; eastern North
America and Australia, colonial-era tales}{A snake that takes its tail in its
mouth and rolls downhill after you, striking with a venomous tail spine. No
snake has a tail stinger and no snake can roll; the story is a tall tale, but a
mechanically specific one, and it has been retold for three hundred years
because it is a good story.}{Mud snakes, which have a pointed tail tip, and
sidewinding or looping defensive motion misread as a wheel.}

\cryptidcard{cryp_mongoliandeathworm}{Mongolian Death Worm}{Desert reptile or
invertebrate; Gobi Desert, Mongolia, reported from the 1920s}{The
\emph{olgoi-khorkhoi}, a red intestine-shaped animal said to spray acid and kill
at a distance. Roy Chapman Andrews recorded the belief in the 1920s and noted
that no Mongolian official he asked had actually seen one --- each had heard of
it from someone else.}{The Tartar sand boa, whose blunt tail and burrowing gait
fit the shape, plus worm lizards of the genus \emph{Amphisbaena}.}

\cryptidcard{cryp_titanoboa}{Titanoboa Survivalists}{Giant snake claim; Amazon
Basin, modern reports}{Claims of anacondas beyond fifty feet, usually supported
by a foreshortened photograph. This card is included because it is the cleanest
demonstration in the book of why measurement matters: the record verified
anaconda is about twenty-three feet, and every metre beyond that has come from
an estimate rather than a tape.}{Green anacondas of genuinely large but
ordinary size, photographed with the head nearest the lens.}

\cryptidcard{cryp_tsuchinoko}{Tsuchinoko}{Venomous serpent; mountains of
Japan, folklore with modern bounties}{A short thick jumping snake, taken
seriously enough that several Japanese towns have posted cash rewards --- one of
them over a hundred thousand US dollars --- and none has ever been claimed. A
standing bounty in a densely populated, thoroughly hiked country is a strong
negative result.}{Japanese mamushi vipers that have swallowed large prey, giving
the swollen wedge-shaped body of the legend.}

% ---------- 6. ANOMALOUS HUMANOIDS & HYBRIDS ----------

\cryptidcard{cryp_doverdemon}{Dover Demon}{Humanoid; Dover, Massachusetts,
21--22 April 1977}{A pale, large-headed, thin-limbed figure with orange eyes,
seen by four teenagers at three separate locations across two nights, each of
whom drew it before comparing notes. The drawings match. Then it stopped, in a
wealthy town twenty miles from Boston, and has never recurred.}{A moose calf or
foal at the edge of headlights, whose leggy silhouette and eyeshine fit --- and
which the witnesses have always rejected.}

\cryptidcard{cryp_flatwoods}{Flatwoods Monster}{Humanoid; Flatwoods, West
Virginia, 12 September 1952}{A ten-foot hovering figure in a spade-shaped cowl,
seen by seven people --- a mother, six children and a dog --- who became
nauseated afterwards, most likely from a mist they walked through. A meteor was
tracked across three states that evening and the site had a barn owl in a
tree.}{A barn owl in a maple, lit from below by flashlights, plus a documented
meteor and irritant mist.}

\cryptidcard{cryp_foukemonster}{Fouke Monster}{Hominid; Fouke, Arkansas,
reports from 1946, wave in 1971}{A red-eyed swamp hominid that scratched at the
Ford house in May 1971, made the wire services, and then made \emph{The Legend
of Boggy Creek} --- a film shot with local residents playing themselves, which
is exactly the mechanism by which a local scare becomes a permanent
cryptid.}{Black bears, with three-toed track casts that plaster experts judged
to have been carved.}

\cryptidcard{cryp_goatman}{Goatman}{Hybrid; Prince George's County, Maryland,
and central Texas, from 1957}{An axe-carrying half-goat man, attached in
Maryland to a story about a federal animal-research escape and in Texas to a
bridge. Both variants are the same story a folklorist can date, and both are
attached to a specific road where teenagers park.}{Local legend-tripping
folklore, with escaped livestock and hoaxers in costume accounting for the
sightings.}

\cryptidcard{cryp_hopkinsville}{Hopkinsville Goblins}{Humanoid; Kelly,
Kentucky, 21--22 August 1955}{Small silver large-eared beings that besieged the
Sutton farmhouse for four hours while the family fired shotguns through the
walls. Police found the shot holes and the terrified family, and no bodies. The
case gave the world the word ``little green men''.}{Two great horned owls
defending a nest, whose height, ear tufts, silver underside and posture on the
ground fit every detail of the description.}

\cryptidcard{cryp_lizardman}{Lizard Man of Scape Ore Swamp}{Reptilian humanoid;
Bishopville, South Carolina, 1988}{A seven-foot green three-fingered thing that
chased a teenager's car and left claw marks in the sheet metal. It is a rare
case with physical damage on record --- and the sheriff's department also
documented the county's insurance claims and the tourist traffic that
followed.}{A black bear standing upright, plus vehicle damage that was
measured, photographed and never conclusively explained.}

\cryptidcard{cryp_lovelandfrogman}{Loveland Frogman}{Amphibian humanoid;
Loveland, Ohio, 1955 and 1972}{Four-foot bipedal frog-like creatures reported
first by a businessman in 1955 and then in 1972 by two police officers on
separate nights, one of whom shot at it. In 1999 the second officer stated
publicly that what he had seen was a large escaped pet iguana, tailless.}{A
large iguana missing its tail --- identified by the original police witness
himself.}

\cryptidcard{cryp_popobawa}{Popobawa}{Shape-shifting humanoid; Zanzibar,
Tanzania, from 1965, panics into the 1990s}{A one-eyed bat humanoid said to
assault sleepers at night, whose appearances cluster tightly around Zanzibari
election years. It belongs on this roster because it shows a cryptid functioning
as politics: the panics are real, mass, dateable social events with real
injuries and real deaths.}{Sleep paralysis, which produces exactly this
experience worldwide, amplified by documented political panic.}

% ---------- 7. FEARSOME CRITTERS (NORTH AMERICAN LUMBERJACK FOLKLORE) ----------

\cryptidcard{cryp_hidebehind}{Hidebehind}{Fearsome critter; northern woods of
the United States, logging-camp lore}{A predator that is always behind the tree
you just passed, and cannot be seen because it hides too well. It is a joke and
an unfalsifiable claim in one sentence --- which makes it a useful teaching tool,
because the reasoning it parodies is used seriously elsewhere in ufology.}{A
logging-camp tall tale, written down by William Cox in 1910.}

\cryptidcard{cryp_hodag}{Hodag}{Fearsome critter; Rhinelander, Wisconsin,
hoaxed 1893}{A horned, spine-backed, green-scaled beast that Eugene Shepard
photographed, then exhibited to paying crowds at county fairs --- until the
Smithsonian announced it was sending investigators, at which point Shepard
admitted the animal was wood, wire and ox hide.}{An admitted hoax, and one of
the best documented in American history --- Rhinelander still uses the Hodag as
its town mascot.}

\cryptidcard{cryp_jackalope}{Jackalope}{Fearsome critter; Douglas, Wyoming,
created 1932}{An antlered jackrabbit invented by the Herrick brothers, who
mounted deer antlers on a hare and sold it to a hotel. Wyoming has since named
Douglas its official home, and the state has considered making it the official
mythical creature. The deep root is real, though: rabbits infected with
Shope papillomavirus grow horn-like keratin tumours.}{A taxidermy invention,
with Shope papillomavirus in wild rabbits producing genuinely horned
animals.}

\cryptidcard{cryp_squonk}{Squonk}{Fearsome critter; hemlock forests of
Pennsylvania, logging-camp lore}{A creature so ashamed of its warty skin that it
weeps constantly, and dissolves entirely into a pool of its own tears when
captured. The story exists to explain why nobody has ever produced one --- the
evidence is designed to destroy itself, which is the shape of a great many
claims in this book.}{Pure folklore, printed by William Cox in 1910; the
self-erasing evidence is the point of the joke.}

\cryptidcard{cryp_slenderman}{Slenderman}{Fearsome critter of the internet era;
created on the Something Awful forums, 10 June 2009}{A faceless figure in a black
suit, impossibly tall, photographed at the edge of a treeline behind playing
children. It has a documented birth certificate: Eric Knudsen posted two doctored
photographs to a Photoshop contest, other users added sightings, and within three
years it had folklore, ritual rules and a mythology nobody had authored. It ends
the roster because it is the control on the control group --- a cryptid whose
origin is known to the hour, that behaved in every other respect exactly like the
fifty-nine entries above it.}{Nothing. An invented character, born 2009, with a named
author --- and believed in anyway.}

\end{multicols}

\rosternarrow
\swampfieldoff


% ---------------------------------------------------------------
\section{Skinwalker Ranch}\label{sec:skinwalker}

This is where the thread ends, and it ends here because Skinwalker Ranch is not one more case.
It is all of them, on one deed, at one address.

A roughly 500-acre property in the Uinta Basin, Utah, near Fort Duchesne, has been the site of
reported anomalies since at least the 1970s: cattle mutilations, orbs, poltergeist activity,
large unidentified animals, disappearing livestock, magnetic and radiation anomalies, and injuries
to investigators.

The history is well documented. The Sherman family bought the ranch in 1994 and reported
extensive phenomena. In 1996 Robert Bigelow, the aerospace entrepreneur, bought the property and
ran the National Institute for Discovery Science on it, staffing it with scientists including
Colm Kelleher and involving journalist George Knapp; their book \emph{Hunt for the Skinwalker}
appeared in 2005. NIDS closed in 2004. Bigelow later held the AAWSAP contract discussed in
Chapter~\ref{ch:government} --- and this connection is genuinely important, because the Defense Intelligence Agency
programme that produced part of the 2017 story was administered by the man who owned the ranch,
and some of AAWSAP's work involved it. Brandon Fugal bought the property in 2016 and the History
channel series began in 2020.

My assessment, and I will be direct.

The research problem is structural. Nearly all reported phenomena come from residents and
investigators who were financially, professionally or emotionally invested. NIDS operated for
eight years with equipment and staff and published essentially nothing in peer-reviewed venues.
The television series generates anomalies weekly on a production schedule, which is the least
plausible feature of the entire story --- genuine anomalous phenomena do not respect a filming
calendar. Access is controlled, and independent researchers have been unable to conduct
unsupervised work.

The environment supplies a lot. The Uinta Basin is a major oil and gas region with substantial
methane, hydrogen sulphide and natural seepage; it hosts high electromagnetic activity from
industrial operations; it is under military airspace, with the Utah Test and Training Range not
far away; and the geology includes uranium-bearing formations, which accounts for background
radiation readings that startle people who have not measured background elsewhere. \hi{Every
radiation or magnetic ``anomaly'' reported without a control measurement is not a measurement.}

And there is one piece of the environment that the show found itself, aired once, and then read
exactly backwards. In Season~2, Episode~8 --- ``Shocking Revelations'' --- Travis Taylor and the
team drove metal stakes into the ground at Homestead~2, fifty to sixty feet apart, wired them to a
12-volt car battery to close a circuit through the soil, and then dropped a 12-volt incandescent
bulb from a Pelican case into the loop to make the current visible. The bulb lit. Roughly a third
of an amp was crossing that stretch of Utah dirt. The moment was presented as evidence that
something extraordinary lies beneath that field. I think it is the single most important
measurement anyone has taken on that property, and the programme drew the one conclusion from it
that cannot be right.

Start with where the current came from, because it did not come from the ground. It came from the
battery. In that circuit the soil was not a source, it was the load --- the resistance the battery
had to push through --- so lighting the lamp measures precisely one thing: how much fifty-odd feet
of this particular dirt, plus two stakes, resists a current you supply yourself. Twelve volts at
about 0.3~amps is a total circuit resistance in the region of forty ohms, and most of that sits in
the few inches of soil gripping each stake rather than in the long run between them, which for
electrodes of that size implies a bulk resistivity in the low tens of ohm-metres. That is wet,
salty ground. \hi{A conductor is not a source. The ranch did not light that bulb; Travis's battery
did, and the dirt merely failed to stop it.}

The distinction matters, because there is a real effect this gets confused with. Two dissimilar
metals pushed into damp soil genuinely do generate current with no battery present at all: the
\term{earth battery}, patented by Alexander Bain in 1841 and used to run telegraph circuits before
dry cells were practical. That is the phenomenon the ranch test would have to have been in order to
mean what the narration implied --- ground producing power of its own --- and it was not, because
there was a car battery in the loop. Either way the destination is the same. The bulb was not
detecting the ranch. It was detecting the dirt.

\actualq{Why did the bulb light --- what is down there powering it?}{Is this dirt an insulator or not? The bulb lit because it is not, and saline shale soil never was.}

And the dirt in that basin is exactly the dirt you would order if you wanted this result. The
Uinta Basin's soils are largely derived from Mancos Shale --- saline and sodic, heavy in sulphates
and sodium --- and the shallow groundwater is saline too, salty enough that USGS geophysical work
in the basin records pore-water resistivities below one ohm-metre in places. Low resistivity is
high conductivity. That is a landscape that carries current.

\hi{Which gives you the whole chain, and it is his chain, not mine to claim: conductive soil means
current flows, current flowing means voltages appear where you did not put them, and moving current
means magnetic fields.} Not in the mesa. In the ground --- the whole property --- and it is natural.
Every electrical oddity the series has built episodes around lives downstream of that one fact:
unexplained voltages on fences and equipment, meters twitching, instruments behaving badly in one
spot and not another, compasses and magnetometers disagreeing with each other. A conductive,
saline soil column is the mechanism. You do not need anything beyond it.

Nor is all of that current natural, which makes it worse for the anomaly reading rather than
better. Rural power distribution in the American West is multigrounded --- the neutral is bonded to
earth repeatedly along the line --- so return current genuinely flows through the soil, and the
resulting \term{neutral-to-earth voltage} (farmers know it as stray voltage, and it is a
well-documented agricultural nuisance) turns up on anything metal and earthed. Add the gas
infrastructure: buried pipelines in an oil and gas basin are protected from corrosion by
\term{cathodic protection}, which works by deliberately pushing direct current into the ground.
So the ranch sits on saline conductive soil, threaded with the return currents of a rural grid and
the protective currents of a pipeline network. Two electrodes and a bulb found all of it.

And you have met this instrument before. Section~\ref{sec:ghosthunting} was about a field whose
entire kit is built around the EMF meter, on the premise that a field reading is a presence ---
so by now the reader knows that in the paranormal trades, electromagnetism \emph{is} the evidence.
\hi{Skinwalker Ranch is that same instrument pointed at ground which genuinely does carry current.}
The ghost hunter's meter spikes because a phone hunts for signal in a basement; the ranch's meters
spike because saline soil, a multigrounded grid and a cathodically protected pipeline network are
all doing exactly what they do. In both cases the detection is real, the field is real, and the
inference drawn from it is the part that fails.

There is one more consequence, and it cuts against the programme's own headline claims. Highly
conductive soil is the classic killer of ground-penetrating radar: conductive, clay-rich, saline
ground attenuates the signal severely and can reduce useful penetration to a fraction of a metre.
The same property that lights the bulb is the property that makes the show's subsurface radar
results --- the mysterious buried structures, the anomalies under the mesa --- least trustworthy.
You cannot invoke the ground's electrical strangeness as evidence and then read its radar returns
as though the ground were ordinary. It is one soil.

Then comes the part the episode treats as its revelation. The instant Travis broke the circuit ---
pulled the connection, killing the current --- Erik Bard's lightning detector fired. On screen this
lands as the property answering back: interfere with the ground and something responds,
immediately, repeatably. \hi{It is repeatable. It is also the single most ordinary thing that can
happen when you open a live circuit.} Interrupting a current does not let it stop politely; the
collapsing magnetic field around the loop dumps its energy into a voltage spike across the break,
which arcs at the contact and radiates a broadband burst of radio noise. That is \term{inductive
kick} --- the same effect that pits relay contacts, the reason engineers solder a snubber across
switched coils, the reason an electric drill makes your radio hiss. And here the loop was a hundred
feet of wire and soil, unshielded, lying on the ground: a serviceable transmitting antenna.
\emph{Lightning} detectors are built to trip on precisely this signature, because that is what
lightning is --- a fast current interruption radiating at low frequency. Bard's instrument was
working perfectly. It detected a spark, a few paces away, made by his colleague's hand.

So the celebrated correlation is real and means nothing anomalous: cause and effect ran from the
team to the instrument, not from the ground to either. Test it the way you would test anything.
Switch the same circuit a hundred times and count the detections; do it in a supermarket car park
in Vernal; put a diode across the break and watch the response die. If the detector fires wherever
you open the loop, the ranch is not in the causal chain. As far as I can establish, none of those
controls were run, or at least none aired.

The same goes for what the later multi-spectral and magnetic surveys reportedly found underneath
Homestead~2: unusual \emph{linear} metallic deposits, described on air as behaving like an
artificial underground grid. Take the linearity seriously, because it is the most diagnostic detail
available, and then ask what makes straight lines underground. Geology mostly does not. People
routinely do: pipelines, well casings, abandoned casing strings and drill pipe in a century-old oil
and gas basin, buried irrigation and water lines, telephone and power runs, homestead plumbing,
fence wire, rebar, dumped machinery. Homestead~2 is not a wilderness site; it is the footprint of a
working farmstead on a producing field. A magnetically detectable straight line under an old ranch
is, on the evidence presented, infrastructure until someone digs it up and shows me otherwise ---
and a metal grid in the soil is also a fine explanation for stray current and misbehaving compasses,
which is to say it belongs in the mechanism, not in the mystery.

What gets skipped, every time, is the boring follow-up. The bulb lit, the detector chirped, and the
programme moved to Tesla coils --- bigger apparatus, better pictures, no controls. That is the
failure mode of the whole enterprise in miniature: the measurements that would have explained the
phenomena were less televisual than the phenomena, so they were not pursued.

\begin{funfact}[The Bulb Test]
\funfactlead{Telegraph engineers were doing this deliberately in the 1840s.} Earth batteries were a
working technology --- Bain's 1841 patent, later ``ground battery'' installations on telegraph
lines --- precisely because ordinary damp soil is a serviceable electrolyte. When two rods in the
ground light a lamp, the historically literate response is not ``what is under this mesa'' but
``how saline is this soil, and what is my control site fifty miles away doing?'' The one difference
worth keeping straight: Bain's ground batteries \emph{generated} the current. At Homestead~2 the
current came out of a car battery, and the soil only carried it.
\end{funfact}

\begin{srcnotes}
\item Bain, A.\ (1841). British patent for the earth battery; see also standard telegraph-era
accounts of ``ground batteries'' as line power.
\item US Geological Survey, Open-File Report 87-394 --- geophysical and pore-water resistivity data
for the Uinta Basin, Utah.
\item Natural Resources Conservation Service soil surveys, Uintah County, Utah --- saline and sodic
soils on Mancos Shale parent material.
\item Doolittle, J.\ A., et al.\ (2007), on soil electrical conductivity and ground-penetrating
radar suitability; US EPA/NRCS GPR soil suitability guidance.
\item USDA and land-grant extension literature on neutral-to-earth (stray) voltage on
multigrounded rural distribution systems.
\item NACE/AMPP cathodic protection standards --- impressed-current protection of buried pipelines.
\item \emph{The Secret of Skinwalker Ranch}, Series~2, Episode~8, ``Shocking Revelations''
(History, 2021) --- the Homestead~2 electrode, battery and bulb test, and the lightning-detector
response on breaking the circuit.
\item Standard electrical-engineering treatments of inductive switching transients and contact arcing
--- inductive kick, snubber and freewheel-diode practice; and the low-frequency broadband emission
that lightning detectors are designed to trigger on.
\end{srcnotes}

And the cultural dimension deserves a great deal more than the passing mention it usually gets,
because the name on the gate is not a piece of colour. It is somebody else's religion, used as
branding, and the details matter.

\figwrap[l]{0.38}{Ute men photographed with a US Army lieutenant near Fort Duchesne, Utah, in 1886, eight years before the reservation was broken up by allotment. The ranch sits on this land.}{https://commons.wikimedia.org/wiki/File:Lt._Styer,_and_Ute_Amigos._Near_Fort_Duchesne,_Utah,1886_-_NARA_-_533076.jpg}{fig:utefortduchesne}
Start with the ground. The Uinta Basin is Ute country --- the Uintah and Ouray Reservation of the
Ute Indian Tribe, whose tribal headquarters are at Fort Duchesne, a couple of miles from the
property (Figure~\ref{fig:utefortduchesne}). The Ute were removed there from Colorado in the
1860s and 1880s; the reservation was then opened to allotment under the Dawes framework in 1905,
which is the mechanism by which land inside a reservation boundary comes to be owned by people who
are not Ute. That is the actual history of the parcel: not a haunted homestead, a dispossession.
The ranch was long known locally as the Sherman place and the anomaly stories circulated in the
Basin before anybody put a word on a sign.

Now the word. \emph{Skinwalker} is an English rendering of \emph{yee naaldlooshii} --- roughly,
by means of it, it goes on all fours --- and it is \emph{Navajo}, not Ute. It refers to a witch:
a person who has deliberately acquired the power to take animal form by committing an act so
monstrous that it inverts them, in most accounts the killing of a close relative. This is not a
monster in the campfire sense. It is a category of human evil, and the reason Navajo people do not
discuss it casually is not superstition about attracting attention --- it is that the subject is
adjacent to accusation. Witchcraft accusation has a documented and lethal history on the Navajo
Nation, and the anthropological literature on it, from Clyde Kluckhohn's \emph{Navaho Witchcraft}
(1944) onwards, reads less like folklore collection than like a study of how a community handles
suspicion. \hi{Talking about skinwalkers in Navajo is closer to talking about paedophiles than to
talking about vampires.} That is the register the branding is borrowed from.

Ute tradition, meanwhile, has its own beings and its own protocols, and they are not
interchangeable with Navajo ones. Ute narrative includes water babies --- small beings associated
with springs and rivers, encountered across the Great Basin, dangerous to the incautious --- and
the bear as an ancestral figure central to the Bear Dance, the oldest Ute ceremony, held each
spring. Sinawav, the creator, and Coyote carry the Ute origin narrative, which explains the
scattering of peoples across the land. Members of the Ute Indian Tribe have made the position
plain in interviews: the phenomena talk is a nuisance, the name is wrong, and being cast as the
local supplier of a curse for a television programme is not a role anybody asked for. Navajo
writers, among them Brenda Norrell and various Din\'e commentators, have objected to the term's use
as entertainment for the same reason a family objects to its surname on a novelty product.

So the naming does three things at once, and each is worth separating. It transplants a Navajo
concept onto Ute land, which is the sort of error that only looks small if you think of Native
people as one thing. It converts a serious moral category into a marketing asset. And it does the
reverse of what it appears to do: the name is not an acknowledgement of indigenous knowledge, it is
a borrowed atmosphere, because none of the indigenous content ever gets consulted about what is
actually happening on that land. The ancient astronaut habit dissected in
Section~\ref{sec:starpeople} is here in a purely commercial form --- \hisoft{take the frightening
part, discard the meaning, keep the brand.}

There is also a straightforward point about who benefits. The Basin's anomaly economy --- the
series, the books, the tours, the merchandise --- routes revenue to a property owner and a
production company. The tribe on whose historical territory the property sits, and whose
neighbouring belief system supplies the title, receives nothing, and its actual concerns in the
Uinta Basin are considerably less telegenic: water rights, oil and gas royalties, jurisdiction,
and the air quality consequences of the same gas field that so usefully explains a number of the
ranch's lights and smells.

What would change my mind is not complicated: continuous instrumented monitoring, with published
calibration and control data, conducted by researchers with no stake in the outcome, with raw
data released. The Hessdalen project in Section~\ref{sec:paranormalskies} has done exactly that in Norway for forty
years on a fraction of the budget. The difference between the two is the difference between
research and content.

Now the reason this section is last.

Run back down the thread and check the inventory. Mutilated cattle, explained by scavengers and
ordinary decomposition (Section~\ref{sec:cattlemutilations}) --- present. Large unidentified
animals that never leave a body, in an era of trail cameras
(Section~\ref{sec:cryptozoology}) --- present. Lights in a valley with a gas field and a
military range in it, which is Hessdalen's question asked without Hessdalen's instruments
(Section~\ref{sec:paranormalskies}) --- present. Radiological and magnetic readings quoted without
a control, on uranium-bearing geology --- present, and that is the
Section~\ref{sec:trumpetsjericho} error in a different medium: a measurement of the ordinary,
sold as a measurement of the strange. Testimony from a small group with financial, professional and
emotional stakes in the outcome, the engine of Section~\ref{sec:montauk} --- present. A borrowed
name supplying age and depth the case cannot supply itself, exactly as the Philadelphia Experiment
borrowed the \emph{Eldridge} and Montauk borrowed a date
(Section~\ref{sec:philadelphiaexperiment}) --- present. A prestige institutional connection, in
Bigelow and the later AAWSAP contract of Chapter~\ref{ch:government}, doing the work that a
peer-reviewed paper would otherwise have to do --- present. And a phenomenon that obliges a
production schedule, which is the crop circle economics of Section~\ref{sec:cropcircles} with
better cameras and a broadcast slot --- present, and decisive.

\hi{Every failure mode this book has spent fourteen chapters isolating is operating simultaneously
on one property, and that is the finding.} Not that the Uinta Basin is uninteresting --- it is a
remarkable place, geologically and historically --- but that Skinwalker Ranch is the field's
methodology assembled in one location and pointed at a camera. It is what you get when investment,
inherited legend, uncontrolled measurement, invested witnesses, institutional borrowing and a
commissioning editor are allowed to compound on the same 500 acres for thirty years.

Which is also why I have not written a conclusion that tells you what to believe. If the thread
from Section~\ref{sec:shadow} onward has done its job, you did not need me for this section.
You had the whole list before you arrived.

\begin{funfact}
The Uinta Basin is one of the most intensively surveyed pieces of ground in the American West ---
not for anomalies, but for hydrocarbons. Thousands of wells, decades of seismic survey, continuous
air quality monitoring, published geology down to the formation. If a 500-acre parcel inside it
had an unexplained magnetic or radiological signature, the people who would have found it first
are the ones who were there for the gas.
\end{funfact}
